\documentclass[11pt]{amsart}

\usepackage{amsmath}
\usepackage{amssymb}
\usepackage{amsthm}
\usepackage{booktabs}
\usepackage{array}
\usepackage{hyperref}

% ---------------------------------------------------------------------------
% Theorem environments
% ---------------------------------------------------------------------------
\newtheorem{theorem}{Theorem}[section]
\newtheorem{proposition}[theorem]{Proposition}
\newtheorem{lemma}[theorem]{Lemma}
\newtheorem{corollary}[theorem]{Corollary}
\theoremstyle{definition}
\newtheorem{definition}[theorem]{Definition}
\theoremstyle{remark}
\newtheorem{remark}[theorem]{Remark}

\numberwithin{equation}{section}
\setlength{\emergencystretch}{2em}

% ---------------------------------------------------------------------------
% Macros
% ---------------------------------------------------------------------------
\newcommand{\NN}{\mathbb{N}}
\newcommand{\ZZ}{\mathbb{Z}}
\newcommand{\QQ}{\mathbb{Q}}
\newcommand{\RR}{\mathbb{R}}
\newcommand{\CC}{\mathbb{C}}
\newcommand{\cN}{\mathcal{N}}
\newcommand{\cE}{\mathcal{E}}
\newcommand{\cS}{\mathcal{S}}
\newcommand{\cP}{\mathcal{P}}
\newcommand{\Acert}{A_{\mathrm{cert}}}
\newcommand{\Abody}{A_{\mathrm{body}}}
\newcommand{\Vcert}{V_{\mathrm{cert}}}
\newcommand{\etafar}{\eta_{\mathrm{far}}}
\newcommand{\xtop}{x_{\mathrm{top}}}
\renewcommand{\Re}{\operatorname{Re}}
\renewcommand{\Im}{\operatorname{Im}}
\newcommand{\eps}{\varepsilon}
\newcommand{\ind}[1]{[\,#1\,]}
\newcommand{\lean}[1]{{\def\_{\textunderscore\allowbreak}\texttt{#1}}}

\begin{document}

\title{One of $\zeta(7),\zeta(9),\dots,\zeta(21)$ is irrational}

\author{Shawn Krisman}

\subjclass[2020]{Primary 11J72; Secondary 11M06, 41A60, 68V20}
\keywords{Riemann zeta function, odd zeta values, irrationality, hypergeometric linear forms,
saddle-point method, formal verification, Lean}

\date{\today}

\begin{abstract}
We prove that at least one of the eight numbers $\zeta(7),\allowbreak\zeta(9),\allowbreak\dots,\allowbreak\zeta(21)$
is irrational. The best previous result of this form that does not involve $\zeta(5)$ was ``one of
$\zeta(7),\dots,\zeta(35)$'', a consequence of the theorem of Lai and Zhou that at least two of
$\zeta(5),\dots,\zeta(35)$ are irrational; Zudilin had earlier proved ``one of $\zeta(7),\dots,\zeta(37)$''. The
proof uses the linear forms $L_n=\sum_{t\ge1}R_n''''(t)$ attached to very-well-poised rational functions $R_n$
of Zudilin's type with quintuple zeros. The fourth derivative and a symmetry keep $\zeta(3)$, $\zeta(5)$ and the
even zeta values out of the forms. The new ingredients are a near-continuous profile of eighteen pole blocks, a
sharper $p$-adic bound at the poles $-k$ with $k<p$, and a saddle-point analysis for the kernel $\pi\cot\pi z$:
four integrations by parts turn the integrand into $R_n(z)$ times an entire trigonometric factor, which fixes the
phase $3\pi$. The common denominators are controlled by an arithmetic integral that we bound with an exact
rational certificate, and the size and nonvanishing of $L_n$ by a certified steepest-descent contour. The whole
proof, including every numerical certificate, has been formalised in Lean~4 with Mathlib and checked by the Lean
kernel, using only Lean's three standard axioms.
\end{abstract}

\maketitle

\tableofcontents

% ===========================================================================
\section{Introduction}\label{sec:intro}
% ===========================================================================

Let $\zeta(s)=\sum_{m\ge1}m^{-s}$ be the Riemann zeta function.

\begin{theorem}\label{thm:main}
At least one of the eight numbers $\zeta(7),\allowbreak \zeta(9),\allowbreak \zeta(11),\allowbreak \dots,\allowbreak \zeta(21)$ is irrational.
\end{theorem}

\subsection{Context}
The even values $\zeta(2k)$ are nonzero rational multiples of $\pi^{2k}$, hence transcendental. Among the odd
values, $\zeta(3)$ is irrational by Ap\'ery's theorem \cite{Apery1979}. Rivoal \cite{Rivoal2000} and Ball and
Rivoal \cite{BallRivoal2001} proved that infinitely many odd values are irrational, and
Fischler, Sprang and Zudilin \cite{FSZ2019} sharpened the count. For short windows, Zudilin
\cite{Zudilin2001} proved that one of $\zeta(5),\zeta(7),\zeta(9),\zeta(11)$ is irrational (see also
\cite[Thm.~3]{Zudilin2004}); Rivoal \cite{Rivoal2002} and Zudilin \cite[Thm.~1]{Zudilin2002} proved it for
the nine numbers $\zeta(5),\dots,\zeta(21)$; and Zudilin gave an elementary proof for
$\zeta(5),\dots,\zeta(25)$ \cite{Zudilin2018}.

A statement ``one of $\zeta(5),\dots,\zeta(N)$ is irrational'' holds as soon as $\zeta(5)$ is irrational, and
then it says nothing about $\zeta(7),\zeta(9),\dots$. The known explicit results about windows that start at
$\zeta(7)$ are the following.
\begin{itemize}
\item Zudilin \cite[Thm.~1]{Zudilin2002} proved that one of $\zeta(7),\zeta(9),\dots,\zeta(37)$ is irrational,
  and more generally \cite[Thm.~2]{Zudilin2002} that for every odd $b\ge1$ one of
  $\zeta(b+2),\zeta(b+4),\dots,\zeta(8b-1)$ is irrational. (Theorem numbers and statements of \cite{Zudilin2002}
  refer to the arXiv version, arXiv:math/0104249, where we checked them.)
\item Lai and Zhou \cite{LaiZhou2022} proved that at least two of $\zeta(5),\zeta(7),\dots,\zeta(35)$ are
  irrational. This implies, with no hypothesis on $\zeta(5)$, that one of $\zeta(7),\dots,\zeta(35)$ is
  irrational: if $\zeta(7),\dots,\zeta(35)$ were all rational, at most one number in the list
  $\zeta(5),\dots,\zeta(35)$ could be irrational.
\item Rivoal and Zudilin \cite{RivoalZudilin2020} proved that at least two of $\zeta(5),\dots,\zeta(69)$ are
  irrational, which gives the same conclusion for $\zeta(7),\dots,\zeta(69)$.
\end{itemize}
So the window $\zeta(7),\dots,\zeta(35)$ was the shortest of this kind that we know of. Theorem~\ref{thm:main}
shortens it to $\zeta(7),\dots,\zeta(21)$, which is Rivoal's window of nine numbers with $\zeta(5)$ removed.

A proof that $\zeta(5)$ is irrational has recently been announced in an unrefereed preprint by A.~Fauzan, with
Lean formalisations (for instance \cite{zeta5}). Theorem~\ref{thm:main} does not use the arithmetic nature of
$\zeta(5)$ or $\zeta(3)$ in any way: neither number occurs in our linear forms (Corollary~\ref{cor:form}).
Theorem~\ref{thm:main} neither implies nor is implied by the result of Lai and Zhou. To our knowledge it is new.
Our literature search covered arXiv and the web, but not MathSciNet or zbMATH; we would be grateful to hear of
prior work.

\subsection{The method}
We follow the approach of Ball and Rivoal \cite{BallRivoal2001} and Zudilin \cite{Zudilin2002,Zudilin2004}
through linear forms with rational coefficients, and the construction is not new: it is the derivative form of
Zudilin's well-poised hypergeometric construction \cite[\S8, (8.2)--(8.4)]{Zudilin2004}, in which
$\frac1{(r-1)!}\sum_{t\ge0}R^{(r-1)}(t)$ is summed; here $r=5$, so the derivative has order $4$. For $n$ in an arithmetic progression we put
\[
  R_n(t)=\frac{(2t+n)\prod_{l=1}^{rn}(t-l)^5\prod_{l=n+1}^{(1+r)n}(t+l)^5}{\prod_{j}\prod_{k=\delta_jn}^{(1-\delta_j)n}(t+k)^{m_j}},
  \qquad L_n=\sum_{t\ge1}R_n''''(t).
\]
Differentiating four times raises every pole order by $4$, so $\zeta(3)$ cannot occur; $\zeta(5)$ drops out
because the residues of $R_n$ add up to $0$; and the symmetry $R_n(-n-t)=-R_n(t)$ removes the even zeta values
(Section~\ref{sec:algebra}). Since all zeros and poles are integers there is no $2$-adic or $3$-adic cost of the
kind that half-integer or third-integer shifts produce. As far as we can tell, Zudilin's theorem for
$\zeta(7),\dots,\zeta(37)$ is of the same type, with all poles in a single block; our single-block instance Z1 of
Section~\ref{sec:general} reaches the same window. The analytic method, a saddle point of the Ball--Rivoal--Zudilin type, is also standard.

What is new is the following.
\begin{enumerate}
\item \emph{A near-continuous pole profile.} The denominator of $R_n$ is a product of eighteen blocks, one at
  each offset $\delta_j=j/80$, $1\le j\le18$, with $r=11/12$. Zudilin's parametrisation allows such profiles; the
  profile used here was found by a numerical search (Section~\ref{sec:search}). Even $S=18$ is useful: the top
  index is then $S+3=21$, so the window $\zeta(7),\dots,\zeta(21)$ gets one block more than an odd $S$ would.
\item \emph{An arithmetic refinement.} If $k<p$, the harmonic sums $H_k^{(j)}$ are $p$-integral, and the
  Taylor coefficients of $R_n$ at the pole $-k$ are $p$-integral unless another pole lies at a distance
  $\equiv0\pmod p$ (Corollary~\ref{cor:A}(b$'$)). So such poles cost nothing at $p$, and the others cost less
  than in the standard estimate (Definition~\ref{def:Ep}). This lowers the arithmetic integral by an amount of
  order $1$ per unit of $n$.
\item \emph{The saddle point for the $\cot$ kernel.} $L_n$ is a contour integral of $R_n''''(z)\pi\cot\pi z$.
  Four integrations by parts move the derivatives onto the kernel, and
  $\sin^5u\cdot\cot''''u=2\cos3u+22\cos u$ (Lemma~\ref{lem:trig}) turns the integrand into $R_n(z)$ times an
  entire function whose leading exponential is $e^{-3\pi iz}$. The resulting phase $\Psi(x)-3\pi ix$ has a
  saddle $x_0$ with $\Re\Omega\le-10.1167485$, which gives both the size of $L_n$ and its nonvanishing for
  infinitely many $n$ (Section~\ref{sec:nonvan}). The elementary bound on the real axis would not suffice here.
\item \emph{Certified numerics.} The arithmetic integral is a sum over $30{,}137$ intervals, evaluated exactly in
  rational arithmetic (Section~\ref{sec:A}). The saddle point, the steepest-descent contour and all bounds along it
  are certified by interval arithmetic.
\item \emph{Formalisation.} Every step, including all computations, has been formalised in Lean~4
  \cite{LeanFour} with Mathlib \cite{Mathlib} (Section~\ref{sec:lean}).
\end{enumerate}
The exponent that must be negative is $Q=\Re\Omega+\Acert<-0.1981$, see \eqref{eq:Q}. Here $\Re\Omega$
controls the size of $L_n$ (Proposition~\ref{prop:size}) and $\Acert$ is the arithmetic integral
(Proposition~\ref{prop:certA}). There is no $2$-adic term.

The same machine gives, with larger margins, the windows $\zeta(7),\dots,\zeta(N)$ for $N=23,25,29,31,33,37$
(Section~\ref{sec:general}); those instances rest on certificates computed in Python and Arb and are not
formalised. It does not reach $N=19$ (Section~\ref{sec:N19}).

\subsection{The formalisation}
The statement of Theorem~\ref{thm:main} and its complete proof, including the prime number theorem, complex
Stirling bounds, the contour-integral representation, the arithmetic certificate and the interval-arithmetic
saddle-point certificate, are formalised in Lean~4 with Mathlib and checked by the Lean kernel. The proof
depends only on Lean's three standard axioms (\lean{propext}, \lean{Classical.choice}, \lean{Quot.sound}). It
contains no \lean{sorry} and no \lean{native\_decide}: every proof, including every computation, is checked by
the kernel. The formal statement uses Mathlib's \lean{riemannZeta} and is described in Section~\ref{sec:lean}.

The prose proof below is self-contained, except that certain finite computations are only described: the exact
value of a rational sum with $30{,}137$ terms (Proposition~\ref{prop:certA}) and the interval-arithmetic checks of
the saddle-point data (Theorem~\ref{thm:saddle}). Both are decidable statements about rational data, proved by
\lean{decide +kernel}. The proof in this paper is the one that is formalised; Section~\ref{sec:lean-diff} notes
where the formal proof is organised differently.

\subsection{Organisation and notation}
Section~\ref{sec:construction} defines the class of rational functions and the instance used for
Theorem~\ref{thm:main}. Section~\ref{sec:algebra} shows that $L_n\in\QQ+\QQ\zeta(7)+\dots+\QQ\zeta(21)$.
Section~\ref{sec:arith} constructs a common denominator $D_n$, and Section~\ref{sec:A} bounds its size through
the arithmetic integral. Section~\ref{sec:nonvan} bounds $|L_n|$ and proves that $L_n\ne0$ for infinitely many
$n$. Section~\ref{sec:proof} combines these results. Section~\ref{sec:general} states the general theorem and the
other instances, Section~\ref{sec:lean} describes the formalisation, and Section~\ref{sec:remarks} contains
remarks on the parameter search, the limits of the method and effectivity.

Throughout, $p$ denotes a prime and $v_p$ the $p$-adic valuation on $\QQ$, with $v_p(0)=+\infty$;
$\ZZ_{(p)}$ is the ring of rationals with $v_p\ge0$. $\lfloor\cdot\rfloor$, $\lceil\cdot\rceil$ and $\{\cdot\}$
are the floor, ceiling and fractional part; $\ind{P}$ is $1$ if $P$ holds and $0$ otherwise; $\#$ denotes
cardinality; intervals $[a,b]$ of integers are written like real intervals. $\theta(y)=\sum_{p\le y}\log p$ is
Chebyshev's function. For a power series $g$ in $\eps$, $[\eps^M]g$ is the coefficient of $\eps^M$. We write
$(i)_4:=i(i+1)(i+2)(i+3)$ and $H_k^{(j)}:=\sum_{m=1}^{k}m^{-j}$ (so $H_0^{(j)}=0$).

% ===========================================================================
\section{The construction}\label{sec:construction}
% ===========================================================================

\subsection{The class of rational functions}
\begin{definition}\label{def:param}
A \emph{parameter set} $\cP=(r,(\delta_j,m_j)_{j=1}^J,q)$ consists of $r\in\QQ_{>0}$, offsets
$\delta_j\in\QQ\cap[0,\frac12)$ with multiplicities $m_j\in\NN_{\ge1}$, and $q\in\NN_{\ge1}$ with $rq\in\ZZ$ and
$\delta_jq\in\ZZ$ for all $j$. We assume $S:=\sum_jm_j\ge2$. The integer $n$ always runs through
$q\NN=\{q,2q,3q,\dots\}$. Put
\begin{gather*}
  a_j:=\delta_jn,\quad b_j:=n-a_j,\quad I_j:=[a_j,b_j]\cap\ZZ,\\
  B_k:=\sum_jm_j\ind{k\in I_j}\ (0\le k\le n),\\
  G:=\sum_jm_j(1-2\delta_j)-10r,\qquad \cN:=(2+r)n+1,
\end{gather*}
and let $N_\cP$ be the largest odd integer $\le S+4$ (so $N_\cP=S+4$ for odd $S$ and $S+3$ for even $S$).
\end{definition}

Each $I_j$ is a nonempty integer interval, symmetric under $k\mapsto n-k$, with $\#I_j=(1-2\delta_j)n+1$; so
$B_{n-k}=B_k$ and $0\le B_k\le S$. The \emph{zero set} is
\[
  \mathcal Z_n:=[1,rn]\cup[-(1+r)n,\,-n-1]\subset\ZZ,
\]
two runs of $rn$ consecutive integers on either side of the pole region $[-n,0]$; each zero will have
multiplicity $5$. For $e\in\{0,1\}$ define
\begin{equation}\label{eq:R}
  R_n(t):=\frac{(2t+n)^e\prod_{l=1}^{rn}(t-l)^5\prod_{l=n+1}^{(1+r)n}(t+l)^5}{\prod_{j}\prod_{k\in I_j}(t+k)^{m_j}}
  =\frac{(2t+n)^e\prod_{z\in\mathcal Z_n}(t-z)^5}{\prod_{k=0}^{n}(t+k)^{B_k}},
\end{equation}
and
\begin{equation}\label{eq:L}
  L_n:=\sum_{t=1}^{\infty}R_n''''(t).
\end{equation}

\begin{lemma}\label{lem:counts}
$\sum_{k=0}^nB_k=\sum_jm_j\#I_j=(G+10r)n+S$ and $\deg R_n=e-Gn-S$. In particular $Gn\in\ZZ$. If $G>0$, then
$\deg R_n\le-2$, so $R_n(t)=O(t^{-2})$, $R_n''''(t)=O(t^{-6})$ as $t\to\infty$, and the series \eqref{eq:L}
converges absolutely.
\end{lemma}

\begin{proof}
$\sum_jm_j((1-2\delta_j)n+1)=(G+10r)n+S$, and the numerator has degree $e+10rn$. Hence
$Gn=e+10rn-S-\deg R_n\in\ZZ$. If $G>0$ then $Gn\ge1$ and $\deg R_n\le1-1-S\le-2$.
\end{proof}

\begin{lemma}[Symmetry]\label{lem:sym}
If $e\equiv1+S+Gn\pmod2$, then $\deg R_n$ is odd and $R_n(-n-t)=-R_n(t)$.
\end{lemma}

\begin{proof}
The first claim is Lemma~\ref{lem:counts}. The map $z\mapsto-n-z$ sends $[1,rn]$ onto $[-n-rn,-n-1]$ and back,
and $-n-t-z=-(t-(-n-z))$; the $e$-factor becomes $2(-n-t)+n=-(2t+n)$; and $-n-t+k=-(t+(n-k))$ with
$B_{n-k}=B_k$. So every linear factor changes sign and the multisets of zeros and poles are preserved, whence
$R_n(-n-t)=(-1)^{\deg R_n}R_n(t)=-R_n(t)$.
\end{proof}

If $qG$ is even, then $Gn$ is even for every $n\in q\NN$, and the condition of Lemma~\ref{lem:sym} fixes
$e\equiv1+S\pmod2$ independently of $n$. This is hypothesis (H6) of Section~\ref{sec:general}. In the language
of hypergeometric series, this symmetry together with the factor $2t+n$ makes $R_n$ \emph{very-well-poised}.

\subsection{The instance Z7}
Theorem~\ref{thm:main} uses the following parameter set, called Z7 (the last column of
Table~\ref{tab:instances}).
\begin{itemize}
\item $r=11/12$ and $q=240$: $n=240N_0$ with $N_0\in\NN_{\ge1}$, so $rn=220N_0$ and $(1+r)n=460N_0$.
\item Eighteen blocks $(\delta_j,m_j)=(j/80,1)$, $1\le j\le18$, so $a_j=3jN_0$.
\item $e=1$.
\end{itemize}

\begin{table}[ht]
\centering
\begin{tabular}{@{}p{0.36\textwidth}p{0.58\textwidth}@{}}
\toprule
quantity & value\\
\midrule
$r$, $q$, $e$ & $11/12$, $240$, $1$\\
blocks $(\delta_j,m_j)$ & $(j/80,1)$ for $j=1,\dots,18$\\
$S=\sum_jm_j$, $N_\cP$ & $18$, $21$\\
$\sum_jm_j(1-2\delta_j)$ & $18-\frac{171}{40}=\frac{549}{40}$\\
$G=\sum_jm_j(1-2\delta_j)-10r$ & $\frac{549}{40}-\frac{55}{6}=\frac{547}{120}$; $qG=1094$\\
zero set $\mathcal Z_n$ & $[1,220N_0]\cup[-460N_0,-240N_0-1]$, each of order $5$\\
$\sum_kB_k$ & $3294N_0+18$\\
$\deg R_n$ & $-1094N_0-17$ (odd)\\
$\cN=(2+r)n+1$ & $700N_0+1$\\
\bottomrule
\end{tabular}
\caption{The parameters of Z7.}\label{tab:blocks}
\end{table}

All blocks have $\delta_j\ge\frac1{80}$, so $B_k=0$ for $k<3N_0$: there is no pole at $0$, and every pole of
$R_n$ lies in $[-n+3N_0,-3N_0]$. The $e$-factor vanishes at $t=-n/2$, which is the central pole ($B_{n/2}=18$);
so that pole has order $17$, and the Taylor data below absorb this automatically ($u_e=0$ in
Lemma~\ref{lem:A}). By Lemma~\ref{lem:sym} and $qG=1094\equiv0$, $S=18$, the choice $e=1$ gives
$R_n(-n-t)=-R_n(t)$ for every $n\in240\NN$.

% ===========================================================================
\section{\texorpdfstring{The linear form lies in $\QQ+\QQ\zeta(7)+\dots+\QQ\zeta(21)$}{The linear form lies in Q + Q zeta(7) + ... + Q zeta(21)}}\label{sec:algebra}
% ===========================================================================

In this section $\cP$ is any parameter set with $G>0$ and $e$ as in Lemma~\ref{lem:sym}.

For $0\le k\le n$ define the power series in $\QQ[[\eps]]$
\begin{equation}\label{eq:fk}
  f_k(\eps):=(n-2k+2\eps)^e\prod_{z\in\mathcal Z_n}(-k-z+\eps)^5\prod_{\substack{0\le k'\le n\\k'\ne k}}(k'-k+\eps)^{-B_{k'}},
\end{equation}
with $(d+\eps)^{-1}=\sum_{M\ge0}(-1)^Md^{-M-1}\eps^M$ for $d\ne0$. Note $-k-z\ne0$, because no zero lies in
the pole region $[-n,0]$. If $t_0:=-k$, then $f_k(\eps)$ is the Taylor expansion of $R_n(t_0+\eps)(t_0+\eps+k)^{B_k}$
at $\eps=0$. We write $u_z:=-k-z$, $u_e:=n-2k$ and $d_{k'}:=k'-k$ when $k$ is fixed.

\begin{lemma}[Partial fractions]\label{lem:pf}
Put $a_{i,k}:=[\eps^{B_k-i}]f_k$ for $0\le k\le n$, $1\le i\le B_k$, and $a_{i,k}:=0$ otherwise. Then
$a_{i,k}\in\QQ$ and, for every $t\in\CC\setminus\{-k:B_k\ge1\}$,
\[
  R_n(t)=\sum_{k=0}^{n}\sum_{i=1}^{S}\frac{a_{i,k}}{(t+k)^i}.
\]
\end{lemma}

\begin{proof}
By Lemma~\ref{lem:counts} the numerator of \eqref{eq:R} has smaller degree than the denominator, so $R_n$ is the
sum of its principal parts at the distinct poles $-k$. The principal part at $t_0=-k$ is
$\sum_{i=1}^{B_k}a_{i,k}(t-t_0)^{-i}$, where $a_{i,k}$ is the coefficient of $(t-t_0)^{B_k-i}$ in the Taylor
expansion of $R_n(t)(t-t_0)^{B_k}$, which is $f_k$. Since $B_k\le S$, all $i\le S$ are covered.
\end{proof}

\begin{lemma}\label{lem:shift}
For integers $j\ge2$ and $k\ge0$, $\sum_{t\ge1}(t+k)^{-j}=\zeta(j)-H_k^{(j)}$.
\end{lemma}

\begin{proof}
Substitute $m=t+k$.
\end{proof}

\begin{proposition}[Expansion]\label{prop:expansion}
Define
\[
  \tau_i:=(i)_4\sum_{k=0}^{n}a_{i,k},\qquad
  \rho_n:=-\sum_{k=0}^{n}\sum_{i=1}^{S}(i)_4\,a_{i,k}\,H_k^{(i+4)} .
\]
Then $\rho_n,\tau_i\in\QQ$ and $L_n=\rho_n+\sum_{i=1}^{S}\tau_i\,\zeta(i+4)$.
\end{proposition}

\begin{proof}
All poles are $\le0$, so for real $t>0$ Lemma~\ref{lem:pf} may be differentiated termwise:
$R_n''''(t)=\sum_{k,i}a_{i,k}(i)_4(t+k)^{-i-4}$. Sum over $t\ge1$; this is a finite sum of absolutely convergent
series, and Lemma~\ref{lem:shift} with $j=i+4\ge5$ gives the claim.
\end{proof}

\begin{lemma}[Symmetry of the coefficients]\label{lem:tau-even}
$a_{i,n-k}=(-1)^{i+1}a_{i,k}$ for all $i,k$. Consequently $\tau_i=0$ for every even $i$.
\end{lemma}

\begin{proof}
By Lemmas~\ref{lem:sym} and~\ref{lem:pf},
\[
  R_n(t)=-R_n(-n-t)=-\sum_{k,i}\frac{a_{i,k}}{(-n-t+k)^i}=\sum_{k,i}\frac{(-1)^{i+1}a_{i,k}}{(t+(n-k))^i},
\]
and uniqueness of the partial-fraction decomposition gives $a_{i,n-k}=(-1)^{i+1}a_{i,k}$. Summing over $k$ gives
$\tau_i=(-1)^{i+1}\tau_i$.
\end{proof}

\begin{lemma}[The residues add up to zero]\label{lem:tau1}
$\sum_ka_{1,k}=0$; hence $\tau_1=0$.
\end{lemma}

\begin{proof}
As $t\to+\infty$, $tR_n(t)\to0$ because $\deg R_n\le-2$, while $t\sum_{k,i}a_{i,k}(t+k)^{-i}\to\sum_ka_{1,k}$.
\end{proof}

\begin{corollary}\label{cor:form}
$L_n=\rho_n+\sum_{i}\tau_i\,\zeta(i+4)$, where $i$ runs over the odd integers with $3\le i\le S$. Hence
$L_n\in\QQ+\QQ\zeta(7)+\QQ\zeta(9)+\dots+\QQ\zeta(N_\cP)$; for Z7,
\[
  L_n=\rho_n+\sum_{j=3}^{10}\tau_{2j-3}\,\zeta(2j+1)\in\QQ+\QQ\zeta(7)+\QQ\zeta(9)+\dots+\QQ\zeta(21).
\]
In particular $\zeta(3)$ and $\zeta(5)$ do not occur.
\end{corollary}

\begin{proof}
In Proposition~\ref{prop:expansion} only $\zeta(i+4)$ with $1\le i\le S$ occurs. Lemma~\ref{lem:tau-even} removes
the even $i$ and Lemma~\ref{lem:tau1} removes $i=1$. For Z7, $i=2j-3$ with $3\le j\le10$.
\end{proof}

% ===========================================================================
\section{Denominators}\label{sec:arith}
% ===========================================================================

In this section $\cP$ is any parameter set with $G>0$, $e$ is as in Lemma~\ref{lem:sym}, and $n\in q\NN$ with
$n>r$. For $0\le k\le n$ we have $0<|u_z|<\cN$ for $z\in\mathcal Z_n$ (indeed $|u_z|\le(1+r)n$),
$|u_e|\le n<\cN$, and $|d_{k'}|\le n<\cN$.

\subsection{Valuations of the Taylor coefficients}

\begin{definition}
Let $p$ be a prime and $\nu,\lambda\in\ZZ$. A power series $g\in\QQ[[\eps]]$ is
\emph{$(\nu,\lambda)$-bounded at $p$} if $v_p([\eps^M]g)\ge \nu-\lambda M$ for all $M\ge0$.
\end{definition}

\begin{lemma}\label{lem:psge}
Let $p$ be a prime and $\lambda\in\ZZ$.
\begin{enumerate}
\item If $g_1$ is $(\nu_1,\lambda)$-bounded and $g_2$ is $(\nu_2,\lambda)$-bounded at $p$, then $g_1g_2$ is
  $(\nu_1+\nu_2,\lambda)$-bounded at $p$.
\item If $u,c\in\QQ$ with $v_p(u)\ge \nu$ and $v_p(c)\ge \nu-\lambda$, then $u+c\eps$ is
  $(\nu,\lambda)$-bounded at $p$.
\item If $d\in\QQ\setminus\{0\}$ with $v_p(d)\le\lambda$ and $\nu\le -v_p(d)$, then $(d+\eps)^{-1}$
  is $(\nu,\lambda)$-bounded at $p$.
\end{enumerate}
\end{lemma}

\begin{proof}
(1) $[\eps^M](g_1g_2)=\sum_{a+b=M}[\eps^a]g_1\,[\eps^b]g_2$, and each term has valuation
$\ge \nu_1-\lambda a+\nu_2-\lambda b$. (2) is immediate (also when $u=0$). (3) The coefficient of $\eps^M$ is
$(-1)^Md^{-M-1}$, of valuation $-v_p(d)-Mv_p(d)\ge \nu-\lambda M$.
\end{proof}

\begin{lemma}[Lemma A]\label{lem:A}
Let $p$ be a prime, $\lambda\ge0$ an integer and $0\le k\le n$. Let $\nu_E\in\ZZ$, $\nu_z\in\ZZ$
$(z\in\mathcal Z_n)$ and $\nu_{k'}\in\ZZ$ $(k'\ne k,\ B_{k'}\ge1)$ satisfy
\begin{itemize}
\item $v_p(u_e)\ge \nu_E$ and $v_p(2)\ge \nu_E-\lambda$ (only if $e=1$; put $\nu_E:=0$ if $e=0$);
\item $v_p(u_z)\ge \nu_z$ and $\nu_z\le\lambda$ for every $z\in\mathcal Z_n$;
\item $v_p(d_{k'})\le\lambda$ and $\nu_{k'}\le -v_p(d_{k'})$ for every $k'\ne k$ with $B_{k'}\ge1$.
\end{itemize}
Put $\nu:=\nu_E+5\sum_{z}\nu_z+\sum_{k'\ne k}B_{k'}\nu_{k'}$. Then $f_k$ is $(\nu,\lambda)$-bounded at $p$, and
$v_p(a_{i,k})\ge \nu-\lambda(B_k-i)$ for $1\le i\le B_k$.
\end{lemma}

\begin{proof}
Apply Lemma~\ref{lem:psge}(2) to $u_e+2\eps$ and to each $u_z+\eps$ (here $v_p(1)=0\ge\nu_z-\lambda$),
Lemma~\ref{lem:psge}(3) to each $(d_{k'}+\eps)^{-1}$, and Lemma~\ref{lem:psge}(1) to the product \eqref{eq:fk}.
\end{proof}

For a prime $p$ and $0\le k\le n$ define
\begin{gather*}
  c_p(k):=5\,\#\{z\in\mathcal Z_n:p\mid k+z\}+e\,\ind{p\mid n-2k},\\
  P_{p,j}(k):=\#\{k'\in I_j\setminus\{k\}:p\mid k'-k\},\qquad P_p(k):=\sum_jm_jP_{p,j}(k),\\
  \nu_p(k):=5\sum_{z\in\mathcal Z_n}v_p(k+z)-\sum_jm_j\sum_{k'\in I_j\setminus\{k\}}v_p(k'-k),
\end{gather*}
and $\lambda_p:=\lfloor\log_p\cN\rfloor$. Note $\sum_{k'\ne k}B_{k'}g(k')=\sum_jm_j\sum_{k'\in I_j\setminus\{k\}}g(k')$
for any $g$.

\begin{corollary}\label{cor:A}
Let $0\le k\le n$ and $1\le i\le B_k$.
\begin{enumerate}
\item[(b$'$)] If $P_p(k)=0$, i.e.\ no pole $-k'\ne-k$ of $R_n$ has $p\mid k'-k$, then $a_{i,k}\in\ZZ_{(p)}$.
\item[(a)] If $p>n$, then $a_{i,k}\in\ZZ_{(p)}$.
\item[(b)] If $p^2>\cN$, then $v_p(a_{i,k})\ge c_p(k)-P_p(k)-(B_k-i)$.
\item[(c)] $v_p(a_{i,k})\ge \nu_p(k)-\lambda_p(B_k-i)$.
\end{enumerate}
\end{corollary}

\begin{proof}
We apply Lemma~\ref{lem:A}; $v_p(u_z)=v_p(k+z)$.
\begin{enumerate}
\item[(b$'$)] Take $\lambda=0$ and all exponents $0$. The hypotheses hold because all $u$'s are integers and
  $v_p(d_{k'})=0$ for every pole $-k'\ne-k$.
\item[(a)] For $k'\ne k$, $0<|k'-k|\le n<p$, so $P_p(k)=0$; apply (b$'$).
\item[(b)] Take $\lambda=1$, $\nu_E=\ind{p\mid u_e}$, $\nu_z=\ind{p\mid u_z}$ and $\nu_{k'}=-\ind{p\mid d_{k'}}$.
  Since $0<|d_{k'}|<\cN<p^2$, $v_p(d_{k'})=\ind{p\mid d_{k'}}\le1$; and $v_p(u_e)\ge\nu_E$ also when $u_e=0$.
  Here $\nu=c_p(k)-P_p(k)$.
\item[(c)] Take $\lambda=\lambda_p$, $\nu_E=0$, $\nu_z=v_p(u_z)$ and $\nu_{k'}=-v_p(d_{k'})$. Since
  $0<|u_z|,|d_{k'}|<\cN<p^{\lambda_p+1}$, both valuations are $\le\lambda_p$. Here $\nu=\nu_p(k)$.\qedhere
\end{enumerate}
\end{proof}

Part (b$'$) is the new ingredient. It needs no condition on $p$, and it is used for $k<p$, where also the
harmonic sums are $p$-integral (Lemma~\ref{lem:B}): the pole $-k$ then costs nothing at $p$ unless another pole
lies at a distance divisible by $p$. (The $e$-factor term $v_p(n-2k)$ is dropped from $\nu_p(k)$ in (c); this only
weakens the bound.)

\begin{lemma}[Lemma B]\label{lem:B}
Let $p$ be a prime, $\lambda\ge0$, $j\ge1$ and $k\ge0$, and suppose $v_p(m)\le\lambda$ for all $1\le m\le k$. Then
$v_p(H_k^{(j)})\ge -j\lambda$. In particular, for $0\le k\le n$: $H_k^{(j)}\in\ZZ_{(p)}$ if $k<p$;
$v_p(H_k^{(j)})\ge -j$ if $p^2>\cN$; and $v_p(H_k^{(j)})\ge-j\lambda_p$ always.
\end{lemma}

\begin{proof}
Each term $m^{-j}$ has valuation $\ge-j\lambda$. For the special cases use $\lambda=0$ (if $k<p$), $\lambda=1$
(as $k\le n<\cN<p^2$) and $\lambda=\lambda_p$ (as $k<\cN$).
\end{proof}

\subsection{The common denominator}

\begin{definition}\label{def:Ep}
For $0\le k\le n$ with $B_k\ge1$ and a prime $p$ put
\[
  e_p(k):=\begin{cases}
  B_k+4-c_p(k)+P_p(k), & k\ge p,\\
  B_k-1-c_p(k)+P_p(k), & k<p \text{ and } P_p(k)\ge1,\\
  0, & k<p \text{ and } P_p(k)=0,
  \end{cases}
\]
and define the integer $E_p$ by
\[
  E_p:=\begin{cases}
  0, & p>n,\\
  \max\bigl(0,\ \max_{k:\,B_k\ge1}e_p(k)\bigr), & p\le n,\ p^2>\cN,\\
  \max_{0\le k\le n}\bigl(-\nu_p(k)\bigr)+(S+4)\lambda_p, & p^2\le \cN
  \end{cases}
\]
(the cases are applied in this order). Put $D_n:=\prod_{p\le n}p^{E_p}\in\QQ_{>0}$.
\end{definition}

Thus $v_p(D_n)=E_p$ for every prime $p$. We do not need $E_p\ge0$ in the third range, so we only assert
$D_n\in\QQ_{>0}$.

\begin{proposition}[Integrality]\label{prop:integral}
$D_n\rho_n\in\ZZ$ and $D_n\tau_i\in\ZZ$ for all $i$.
\end{proposition}

\begin{proof}
It suffices to show $v_p(\tau_i)\ge -E_p$ and $v_p(\rho_n)\ge-E_p$ for every prime $p$. Since $(i)_4\in\ZZ$, it
suffices to bound $v_p(a_{i,k})$ and $v_p(a_{i,k}H_k^{(i+4)})$ from below by $-E_p$, for $0\le k\le n$ and
$1\le i\le B_k$ (the other $a_{i,k}$ vanish).
\begin{itemize}
\item $p>n$: $a_{i,k}\in\ZZ_{(p)}$ by Corollary~\ref{cor:A}(a), and $H_k^{(i+4)}\in\ZZ_{(p)}$ by
  Lemma~\ref{lem:B} since $k\le n<p$.
\item $p\le n$, $p^2>\cN$, $k\ge p$: by Corollary~\ref{cor:A}(b) and Lemma~\ref{lem:B},
  \[
    v_p\bigl(a_{i,k}H_k^{(i+4)}\bigr)\ge c_p(k)-P_p(k)-(B_k-i)-(i+4)=-e_p(k)\ge-E_p,
  \]
  and $v_p(a_{i,k})\ge c_p(k)-P_p(k)-B_k+1\ge-e_p(k)$.
\item $p\le n$, $p^2>\cN$, $k<p$: $H_k^{(i+4)}\in\ZZ_{(p)}$ by Lemma~\ref{lem:B}. If $P_p(k)\ge1$,
  Corollary~\ref{cor:A}(b) gives $v_p(a_{i,k})\ge c_p(k)-P_p(k)-B_k+1=-e_p(k)$. If $P_p(k)=0$,
  Corollary~\ref{cor:A}(b$'$) gives $v_p(a_{i,k})\ge0=-e_p(k)$.
\item $p^2\le\cN$: by Corollary~\ref{cor:A}(c) and Lemma~\ref{lem:B}, since $B_k\le S$,
  \[
    v_p\bigl(a_{i,k}H_k^{(i+4)}\bigr)\ge\nu_p(k)-\lambda_p(B_k-i)-\lambda_p(i+4)\ge\nu_p(k)-(S+4)\lambda_p\ge-E_p,
  \]
  and the same bound holds without the factor $H_k^{(i+4)}$.\qedhere
\end{itemize}
\end{proof}

\begin{corollary}\label{cor:intform}
$D_nL_n=D_n\rho_n+\sum_{i}(D_n\tau_i)\,\zeta(i+4)$ ($i$ odd, $3\le i\le S$) is a linear form in
$1,\zeta(7),\dots,\zeta(N_\cP)$ with integer coefficients.
\end{corollary}

\subsection{The small primes}

Let $s_p(N)$ denote the sum of the base-$p$ digits of $N\ge0$, and $\ell_p(N):=\lfloor\log_pN\rfloor$ for
$N\ge1$, $\ell_p(0):=0$.

\begin{lemma}[Legendre-type bounds]\label{lem:legendre}
Let $p$ be a prime.
\begin{enumerate}
\item For $N\ge0$: $(p-1)v_p(N!)=N-s_p(N)$ and $0\le s_p(N)\le(p-1)(\ell_p(N)+1)$.
\item For integers $1\le A\le B+1$:
  \begin{align*}
    (p-1)\sum_{d=A}^{B}v_p(d)&\ge(B-A+1)-(p-1)(\ell_p(B)+1),\\
    (p-1)\sum_{d=A}^{B}v_p(d)&\le(B-A+1)+(p-1)(\ell_p(A-1)+1).
  \end{align*}
\end{enumerate}
\end{lemma}

\begin{proof}
(1) is Legendre's formula, and each of the at most $\ell_p(N)+1$ digits is $\le p-1$. (2) The sum equals
$v_p(B!)-v_p((A-1)!)$, so $(p-1)$ times it is $(B-A+1)-s_p(B)+s_p(A-1)$; use (1).
\end{proof}

\begin{lemma}\label{lem:polesum}
Let $p$ be a prime, $0\le k\le n$ and $1\le j\le J$. Then
$(p-1)\sum_{k'\in I_j\setminus\{k\}}v_p(k'-k)\le \#I_j+2(p-1)(\lambda_p+1)$.
\end{lemma}

\begin{proof}
The numbers $|k'-k|$ with $k'\in I_j$, $k'>k$, form an interval $[A_1,B_1]$ of positive integers, and those with
$k'<k$ form another interval $[A_2,B_2]$; their total length is at most $\#I_j$, and $A_i-1\le n<\cN$. Apply the
upper bound of Lemma~\ref{lem:legendre}(2) to each (an empty interval contributes $0$).
\end{proof}

\begin{proposition}\label{prop:smallprimes}
Let $p$ be a prime with $p^2\le\cN$. Then $p\le n$ and
\begin{equation}\label{eq:smallp}
  E_p-G\,v_p(n!)\le(4S+14+G)(\lambda_p+1).
\end{equation}
Consequently
\[
  \sum_{p^2\le\cN}\bigl(E_p-G\,v_p(n!)\bigr)\log p\le2(4S+14+G)\lfloor\sqrt\cN\rfloor\log\cN .
\]
For Z7 the constants are $4S+14+G\le91$ and $2(4S+14+G)\le182$.
\end{proposition}

\begin{proof}
Since $n>r$, $\cN=(2+r)n+1<(n+1)^2$, so $p\le n$. Fix $k$. As $z$ runs through $\mathcal Z_n$, $k+z$ runs
through the two intervals $[k+1,k+rn]$ and $[k-(1+r)n,k-n-1]$, each consisting of $rn$ consecutive nonzero
integers of absolute value $<\cN$. By the lower bound of Lemma~\ref{lem:legendre}(2), applied to the absolute
values,
\[
  (p-1)\cdot5\sum_{z\in\mathcal Z_n}v_p(k+z)\ge10rn-10(p-1)(\lambda_p+1).
\]
By Lemma~\ref{lem:polesum} and Lemma~\ref{lem:counts},
\[
  (p-1)\sum_jm_j\sum_{k'\in I_j\setminus\{k\}}v_p(k'-k)\le(G+10r)n+S+2S(p-1)(\lambda_p+1).
\]
Hence $(p-1)(-\nu_p(k))\le Gn+S+(2S+10)(p-1)(\lambda_p+1)$. By Lemma~\ref{lem:legendre}(1) and
$\ell_p(n)\le\lambda_p$, $(p-1)v_p(n!)\ge n-(p-1)(\lambda_p+1)$. Therefore
\begin{multline*}
  (p-1)\bigl(E_p-Gv_p(n!)\bigr)\le S+(2S+10)(p-1)(\lambda_p+1)\\+(S+4)(p-1)\lambda_p+G(p-1)(\lambda_p+1),
\end{multline*}
and dividing by $p-1\ge1$ gives \eqref{eq:smallp}. For the sum, each term is at most
$(4S+14+G)(\lambda_p\log p+\log p)\le2(4S+14+G)\log\cN$, since $p^{\lambda_p}\le\cN$ and $p\le\cN$, and there are at
most $\lfloor\sqrt\cN\rfloor$ primes in the range. For Z7, $4\cdot18+14+\frac{547}{120}<91$.
\end{proof}

There is no exceptional prime here: $p=2$ and $p=3$ are covered by the same count, because all zeros and poles
of $R_n$ are integers. (For the $\beta$-analogue \cite{BetaPaper}, with half-integer poles, the prime $2$ costs
a term proportional to $n$.)

\begin{proposition}\label{prop:decomp}
For every $n\in q\NN$ with $n>r$,
\[
  \log D_n-G\log n!\le2(4S+14+G)\lfloor\sqrt\cN\rfloor\log\cN
  +\sum_{\substack{p\le n\\p^2>\cN}}\bigl(E_p-G\lfloor n/p\rfloor\bigr)\log p .
\]
\end{proposition}

\begin{proof}
By definition and Legendre's formula, $\log D_n-G\log n!=\sum_{p\le n}(E_p-Gv_p(n!))\log p$. Split the sum into the
primes with $p^2\le\cN$ (Proposition~\ref{prop:smallprimes}) and those with $p^2>\cN$, for which $p^2>\cN>n$ and so
$v_p(n!)=\lfloor n/p\rfloor$.
\end{proof}

% ===========================================================================
\section{The arithmetic integral}\label{sec:A}
% ===========================================================================

In this section we bound the last sum in Proposition~\ref{prop:decomp} by $n(\Acert+o(1))$. The definitions
are for a general parameter set; the certificate is for Z7, where
\begin{equation}\label{eq:Acert}
  \Acert:=\frac{11757994664484187722611840743096153}{1185451381619680151243401356902400}
  =9.9185802528\ldots
\end{equation}

\subsection{The step function}
\begin{definition}\label{def:f}
For real $x\ge1$ and $0\le y\le x$ put
\begin{align*}
  \cS(x,y)&:=\sum_jm_j\,\ind{\delta_jx\le y\le(1-\delta_j)x},\\
  P(x,y)&:=\sum_jm_j\,\#\{m\in\ZZ\setminus\{0\}:\ \delta_jx\le y+m\le(1-\delta_j)x\},\\
  c(x,y)&:=5\bigl(\lfloor y+rx\rfloor-\lfloor y\rfloor\bigr)+5\bigl(\lceil y-x\rceil-\lceil y-(1+r)x\rceil\bigr),
\end{align*}
\[
  b(x,y):=\begin{cases}
  \cS+4-c+P, & y\ge1,\\
  \cS-1-c+P, & y<1\text{ and }P\ge1,\\
  0, & y<1\text{ and }P=0,
  \end{cases}
\]
\[
  g(x,y):=\begin{cases}b(x,y),&\cS(x,y)>0,\\0,&\cS(x,y)=0,\end{cases}
\]
(arguments $(x,y)$ omitted on the right), and
\[
  e^*(x):=\max\Bigl(0,\max_{0\le y\le x}g(x,y)\Bigr)-G\lfloor x\rfloor .
\]
\end{definition}
The maximum exists because $g(x,\cdot)$ takes finitely many values on $[0,x]$. For Z7 no block contains $y=0$,
so $g(x,0)=0$ and the outer $\max(0,\cdot)$ is redundant; the formalisation defines $e^*$ without it. The
formalisation also writes the pole count of block $j$ by cases: $\lfloor y-\delta_jx\rfloor+\lfloor(1-\delta_j)x-y\rfloor$
in the block, and $\lfloor(1-\delta_j)x-y\rfloor-\lfloor\delta_jx-y\rfloor+\ind{\delta_jx-y\in\ZZ}$ (resp.\ the mirror
expression) to the left (resp.\ right) of it; by $\lceil a\rceil=\lfloor a\rfloor+1-\ind{a\in\ZZ}$ this is the
count above.

\begin{proposition}[Pointwise bound]\label{prop:pointwise}
Let $p\le n$ be a prime. Then $\max\bigl(0,\max_{k:B_k\ge1}e_p(k)\bigr)-G\lfloor n/p\rfloor\le e^*(n/p)$. In
particular $E_p-G\lfloor n/p\rfloor\le e^*(n/p)$ whenever $p^2>\cN$.
\end{proposition}

\begin{proof}
Let $0\le k\le n$ with $B_k\ge1$, and put $x:=n/p\ge1$ and $y:=k/p\in[0,x]$. Since $a_j/p=\delta_jx$ and
$b_j/p=(1-\delta_j)x$, we have $B_k=\cS(x,y)$, and, writing $k'=k+mp$,
$P_{p,j}(k)=\#\{m\ne0:\delta_jx\le y+m\le(1-\delta_j)x\}$, so $P_p(k)=P(x,y)$. The zeros $z\in[1,rn]$ with
$p\mid k+z$ correspond to the multiples of $p$ in $[k+1,k+rn]$, of which there are
$\lfloor(k+rn)/p\rfloor-\lfloor k/p\rfloor=\lfloor y+rx\rfloor-\lfloor y\rfloor$. The zeros $z=-l$,
$l\in[n+1,(1+r)n]$, with $p\mid l-k$ correspond to the multiples of $p$ in $[n+1-k,(1+r)n-k]$, of which there are
$\lfloor(1+r)x-y\rfloor-\lfloor x-y\rfloor=\lceil y-x\rceil-\lceil y-(1+r)x\rceil$. Hence
$c_p(k)=c(x,y)+e\ind{p\mid n-2k}\ge c(x,y)$. Finally $k\ge p$ if and only if $y\ge1$. Comparing the three cases
of Definition~\ref{def:Ep} with those of $b$, we get $e_p(k)\le b(x,y)=g(x,y)$. Take the maximum over $k$ and
subtract $G\lfloor x\rfloor$.
\end{proof}

\begin{lemma}[Uniform bound]\label{lem:tail}
For $x\ge1$, $e^*(x)\le S+14+G\{x\}$.
\end{lemma}

\begin{proof}
For each block, $\ind{\delta_jx\le y\le(1-\delta_j)x}+\#\{m\ne0:\dots\}=\#\{m\in\ZZ:\delta_jx\le y+m\le(1-\delta_j)x\}
\le(1-2\delta_j)x+1$. Each of the two zero counts in $c$ is at least $rx-1$, so $c\ge10rx-10$. Hence, if $y\ge1$,
$b=\sum_jm_j(\ind{\cdot}+\#\{\cdot\})+4-c\le S+(G+10r)x+4-10rx+10=S+14+Gx$, and the other two cases of $b$
are smaller or $0$. So $\max(0,\max_yg)\le S+14+Gx$, and $e^*(x)\le S+14+Gx-G\lfloor x\rfloor$.
\end{proof}

For Z7 this reads $e^*(x)\le32+G\{x\}$.

\subsection{Step structure}
Writing $\#\{m\in\ZZ:\delta x\le y+m\le(1-\delta)x\}=1+\lfloor y-\delta x\rfloor-\lceil y-(1-\delta)x\rceil$,
$\ind{\delta x\le y\le(1-\delta)x}=\ind{\lfloor y-\delta x\rfloor\ge0}\cdot\ind{\lceil y-(1-\delta)x\rceil\le0}$
and $\ind{y\ge1}=\ind{\lfloor y\rfloor\ge1}$, we see that $g(x,y)$ depends on $y$ only through the numbers
$\lfloor y-s x\rfloor$ and $\lceil y-sx\rceil$ for the slopes
\[
  s\in\mathcal T:=\{0,\,-r,\,1,\,1+r\}\cup\{\delta_j,\,1-\delta_j:1\le j\le J\}.
\]
(The slopes $0$ and $1$ also carry the constraint $0\le y\le x$.) For Z7, $\mathcal T=\{\sigma/240\}$ with
$\sigma\in\{0,-220,240,460\}\cup\{3j,240-3j:1\le j\le18\}$. Let $X:=400$. The \emph{event set} $\cE\subset[1,X]$
is the union of
\begin{itemize}
\item the points $1+j/20$ $(0\le j<40)$ and $3+j/2$ $(0\le j\le 2(X-3))$, which include every integer in
  $[1,X]$;
\item all $x\in(1,X)$ of the form $x=m/(s_1-s_2)$ with $s_1,s_2\in\mathcal T$, $s_1>s_2$, and $m\in\ZZ$.
\end{itemize}
This is a finite set of rationals, $1=x_0<x_1<\dots<x_L=X$. For Z7 it has $30{,}138$ elements, so
$L=30{,}137$.

\begin{lemma}\label{lem:step}
$e^*$ is constant on each open interval $(x_{\ell-1},x_\ell)$.
\end{lemma}

\begin{proof}
Let $J=(x_{\ell-1},x_\ell)$ and $x,x'\in J$. Consider the discrete set $C(x):=\{sx+m:s\in\mathcal T,\ m\in\ZZ\}$. For
two of its elements the difference $(s-s')x+(m-m')$ is an affine function of $x$; if it is not identically zero,
its only zero is an event point or lies outside $(1,X)$, so it does not vanish on $J$. Hence $sx+m\mapsto sx'+m$
is a well-defined, strictly increasing bijection $C(x)\to C(x')$, which extends to an increasing bijection
$\pi:\RR\to\RR$. For every $s\in\mathcal T$ and every $y$, $sx+m\le y\iff sx'+m\le\pi(y)$, so
$\lfloor\pi(y)-sx'\rfloor=\lfloor y-sx\rfloor$, and likewise for ceilings. Moreover $\pi(0)=0$ and $\pi(x)=x'$
(slopes $0$ and $1$), so $\pi$ maps $[0,x]$ onto $[0,x']$, and $\lfloor x\rfloor=\lfloor x'\rfloor$ because $\cE$
contains every integer of $[1,X]$. Therefore $g(x',\pi(y))=g(x,y)$ for all $y$, and $e^*(x')=e^*(x)$.
\end{proof}

The proof also shows how $e^*$ is computed exactly at a rational point $\xi$: $g(\xi,\cdot)$ is constant between
consecutive points of $C(\xi)$, so the maximum over $[0,\xi]$ is attained at a point of $C(\xi)$ or as a right
limit there. In the formalisation, with $\xi=P/Q$ and $M=240Q$, all points of $C(\xi)$ lie in $\frac1M\ZZ$, and
$e^*(\xi)$ is a maximum over the residues modulo $M$ of the scaled offsets of one closed-form candidate each; that
this maximum equals $\max_{0\le y\le\xi}g(\xi,y)$, and is attained, is proved in general.

\subsection{The certificate}

\begin{proposition}[Certified computation, Z7]\label{prop:certA}
Let $\xi_\ell:=(x_{\ell-1}+x_\ell)/2$ and $\Abody:=\sum_{\ell=1}^{L}e^*(\xi_\ell)\bigl(\frac1{x_{\ell-1}}-\frac1{x_\ell}\bigr)$.
Then
\[
  \sum_{\ell=1}^{L}\Bigl\lceil 2^{96}e^*(\xi_\ell)\Bigl(\frac1{x_{\ell-1}}-\frac1{x_\ell}\Bigr)\Bigr\rceil
  =U:=778841997893779320302794447562 .
\]
Consequently $\Abody\le U/2^{96}=9.8303680\ldots$ and
\[
  \frac{U}{2^{96}}+\frac{33}{400}+\frac{G}{2\cdot 399}=\Acert ,
\]
with $\Acert$ as in \eqref{eq:Acert}. (Note $33/400+G/798=10559/119700$, and $33=S+15$.)
\end{proposition}

\begin{proof}
This is a finite computation with explicit rationals: enumerate $\cE$, evaluate $e^*(\xi_\ell)$ exactly, round each
term up to a multiple of $2^{-96}$, and add. The last identity is exact rational arithmetic. The computation has
been carried out by an exact-rational Python implementation and by the Lean kernel, with the same result; see
Section~\ref{sec:lean-cert}.
\end{proof}

\subsection{From primes to the integral}

\begin{lemma}[Tail]\label{lem:pnt-tail}
Let $X\ge2$ be an integer, $C,G'\ge0$, and $h:[X,\infty)\to\RR$ with $h(y)\le C+G'\{y\}$. For every
$\eps>0$, for all sufficiently large $n$,
\[
  \sum_{p\le n/X}h(n/p)\log p\le n\Bigl(\frac CX+\frac{G'}{2(X-1)}\Bigr)+\eps n .
\]
\end{lemma}

\begin{proof}
Choose an integer $K_0\ge1$ with $G'/(2K_0(X-1))\le\eps/4$ and an integer $M>X$ with
$(C+G')\log4/M\le\eps/4$. The primes $p\le n/M$ contribute at most $(C+G')\theta(n/M)\le(C+G')(\log4)\,n/M\le\eps n/4$,
by Chebyshev's bound $\theta(y)\le y\log4$. Cut $[X,M)$ into the $K_0(M-X)$ cells
$[\,m+\frac{i-1}{K_0},\,m+\frac i{K_0})$, $X\le m<M$, $1\le i\le K_0$. On such a cell $h\le C+G'i/K_0$. The primes
with $n/p$ in the cell $[a,b)$ are those in $(n/b,n/a]$, and by the prime number theorem
$\theta(n/a)-\theta(n/b)=n(1/a-1/b)+o(n)$. Since there are finitely many cells, for large $n$ the primes
$n/M<p\le n/X$ contribute at most
\[
  n\sum_{m=X}^{M-1}\sum_{i=1}^{K_0}\Bigl(C+\frac{G'i}{K_0}\Bigr)
  \Bigl(\frac1{m+\frac{i-1}{K_0}}-\frac1{m+\frac i{K_0}}\Bigr)+\frac{\eps n}4 .
\]
The $C$-part telescopes to $C(1/X-1/M)\le C/X$. In the $G'$-part each bracket is at most $1/(K_0m^2)$, so the inner
sum is at most $G'(K_0+1)/(2K_0m^2)$; since $\sum_{m\ge X}m^{-2}\le1/(X-1)$, the $G'$-part is at most
$\frac{G'}{2(X-1)}+\frac\eps4$.
\end{proof}

\begin{proposition}\label{prop:lemma5}
For Z7 there are $\eps_n\to0$ such that, for all large $n\in240\NN$,
\[
  \sum_{\substack{p\le n\\p^2>\cN}}\bigl(E_p-G\lfloor n/p\rfloor\bigr)\log p\le n(\Acert+\eps_n).
\]
\end{proposition}

\begin{proof}
Fix $\eps>0$ and define $h$ on $[1,\infty)$ by $h(x):=e^*(x)$ for $x<400$ and $h(x):=\max(e^*(x),0)$ for
$x\ge400$. Let $n$ be large; in particular $n/\sqrt\cN\ge400$, which holds as soon as $N_0\ge400^2$.
\begin{enumerate}
\item By Proposition~\ref{prop:pointwise}, each term of the sum is at most $h(n/p)\log p$. Every prime $p\le n$
  not in the sum has $p^2\le\cN$, hence $n/p\ge400$ and $h(n/p)\ge0$. So the sum is at most
  $\sum_{p\le n}h(n/p)\log p$.
\item \emph{Primes $n/400<p\le n$}, i.e.\ $n/p\in[1,400)$. Let $v_\ell:=e^*(\xi_\ell)$. By Lemma~\ref{lem:step},
  $h(n/p)=v_\ell$ when $n/p\in(x_{\ell-1},x_\ell)$, and by the prime number theorem
  $\sum_{x_{\ell-1}<n/p<x_\ell}\log p=n(1/x_{\ell-1}-1/x_\ell)+o(n)$ for each of the finitely many $\ell$. At most
  $L+1$ primes have $n/p\in\cE$, and each contributes $O(\log n)$, since $|e^*|\le S+14+400G$ on $[1,400]$ by
  Lemma~\ref{lem:tail}. Hence this part is at most $n\Abody+\eps n/2$ for large $n$.
\item \emph{Primes $p\le n/400$.} By Lemma~\ref{lem:tail}, $h(x)\le32+G\{x\}\le33+G\{x\}$ for $x\ge400$ (the
  certificate uses $33=S+15$). Lemma~\ref{lem:pnt-tail} with $X=400$, $C=33$, $G'=G$ bounds this part by
  $n(33/400+G/798)+\eps n/2$.
\end{enumerate}
Adding up and using Proposition~\ref{prop:certA}, the sum is at most $n(\Acert+\eps)$ for large $n$.
\end{proof}

The same argument works for any parameter set, with $C=S+15$ and with $\Acert$ the certified bound for that set.

\begin{corollary}\label{cor:logD}
For Z7 there are $\eps'_n\to0$ with $\log D_n\le G\log n!+n(\Acert+\eps'_n)$ for all large $n\in240\NN$.
\end{corollary}

\begin{proof}
Combine Propositions~\ref{prop:decomp} and~\ref{prop:lemma5}; the term $182\lfloor\sqrt\cN\rfloor\log\cN$ is
$O(\sqrt n\log n)=o(n)$.
\end{proof}

% ===========================================================================
\section{\texorpdfstring{Size and nonvanishing of $L_n$}{Size and nonvanishing of L\_n}}\label{sec:nonvan}
% ===========================================================================

We extend $R_n$ to the rational function of $z\in\CC$ given by \eqref{eq:R}. It has real coefficients, and its
poles are the points $-k$ with $B_k\ge1$, all $\le0$ (for Z7, all $\le-3N_0$). Put $k(z):=\pi\cot(\pi z)$ and
$k_j:=k^{(j)}$.

\subsection{An integral representation}

\begin{lemma}\label{lem:vanish}
$R_n''''(t)=0$ for every integer $t\in[1,rn]$.
\end{lemma}

\begin{proof}
Near such $t$, $R_n(s)=(s-t)^5g(s)$ with $g$ holomorphic at $t$, so all derivatives of order $\le4$ vanish at $t$.
\end{proof}

\begin{lemma}\label{lem:decay}
Let $c>\max\{-k:B_k\ge1\}$ and $j\ge0$. There is $C=C(n,c,j)$ with $|R_n^{(j)}(z)|\le C(1+|z|^2)^{-1}$ whenever
$\Re z\ge c$. Moreover $R_n^{(j)}(\bar z)=\overline{R_n^{(j)}(z)}$.
\end{lemma}

\begin{proof}
By Lemma~\ref{lem:counts}, $|R_n(w)|\le C_1|w|^{-2}$ for $|w|\ge R_1$. Cauchy's estimate on circles of radius
$\frac12$ gives $|R_n^{(j)}(z)|\le C_1'|z|^{-2}$ for $|z|\ge2R_1+2$, and on the compact set
$\{\Re z\ge c,\ |z|\le2R_1+2\}$, which contains no pole, $R_n^{(j)}$ is bounded. The symmetry holds because $R_n$ has
real coefficients.
\end{proof}

\begin{lemma}[The kernel]\label{lem:trig}
With $\mathsf C:=\cot(\pi z)$ and $\mathsf S:=1/\sin(\pi z)$,
\begin{gather*}
  k_1=-\pi^2\mathsf S^2,\quad k_2=2\pi^3\mathsf C\mathsf S^2,\\
  k_3=-2\pi^4(\mathsf S^4+2\mathsf C^2\mathsf S^2),\quad k_4=8\pi^5(\mathsf C^3\mathsf S^2+2\mathsf C\mathsf S^4).
\end{gather*}
Consequently $\sin^5(\pi z)\,k_4(z)=\pi^5T(z)$ with the entire function
\begin{align*}
  T(z)&:=8\cos(\pi z)\bigl(\cos^2(\pi z)+2\bigr)=2\cos3\pi z+22\cos\pi z\\
  &=e^{-3\pi iz}\bigl(1+11\mathbf q+11\mathbf q^2+\mathbf q^3\bigr),\qquad \mathbf q:=e^{2\pi iz},
\end{align*}
and $|T(z)|\le24e^{3\pi\Im z}$ for $\Im z\ge0$. The kernel $k$ has a simple pole with residue $1$ at every
integer; $k_0,\dots,k_4$ are bounded on every vertical line $\Re z=c\notin\ZZ$ and on $\{|\Im z|\ge1\}$.
\end{lemma}

\begin{proof}
Use $\mathsf C'=-\pi\mathsf S^2$ and $\mathsf S'=-\pi\mathsf C\mathsf S$ and differentiate. Then
$\sin^5(\pi z)k_4=8\pi^5\cos(\pi z)(\cos^2(\pi z)+2)$, and $8\cos^3u+16\cos u=2\cos3u+22\cos u$. Expanding the
cosines in exponentials gives the $\mathbf q$-form, and $|\mathbf q|\le1$ for $\Im z\ge0$. The last statements
are standard.
\end{proof}

\begin{proposition}\label{prop:contour}
Put $G_n(z):=R_n(z)k_4(z)$. For every real $c\in(0,rn+1)\setminus\ZZ$,
\begin{align*}
  L_n&=-\frac1{2\pi i}\int_{c-i\infty}^{c+i\infty}R_n''''(z)\,k(z)\,dz
     =-\frac1{2\pi i}\int_{c-i\infty}^{c+i\infty}G_n(z)\,dz\\
     &=-\frac1\pi\,\Re\int_0^\infty G_n(c+iy)\,dy .
\end{align*}
\end{proposition}

\begin{proof}
\emph{First equality.} Let $N>c$ be an integer and $Y>1$, and integrate $R_n''''k$ counterclockwise around the
rectangle with vertices $c\pm iY$, $N+\frac12\pm iY$. Inside, the integrand is meromorphic with simple poles at
the integers $t\in(c,N]$, with residue $R_n''''(t)$. On the horizontal sides $k$ is bounded and
$|R_n''''(z)|\le C(1+|z|^2)^{-1}\le CY^{-2}$, so these sides contribute $O(N/Y^2)$, which tends to $0$ as
$Y\to\infty$. On $\Re z=N+\frac12$, $|k|\le\pi$ and $|R_n''''|\le C(1+|z|^2)^{-1}$, so that side tends to $0$ as
$N\to\infty$. On $\Re z=c$ the integral converges absolutely. Hence
$-\int_{c-i\infty}^{c+i\infty}R_n''''k\,dz=2\pi i\sum_{t>c}R_n''''(t)=2\pi i\,L_n$, where the last equality
uses Lemma~\ref{lem:vanish}: every integer $1\le t<c$ satisfies $t\le rn$.

\emph{Second equality.} On the line $\Re z=c$ put
$W:=R_n'''k_0-R_n''k_1+R_n'k_2-R_nk_3$. Then $W'=R_n''''k_0-R_nk_4$ (the other terms cancel in pairs). The $k_j$
are bounded on the line and $R_n^{(j)}(c+iy)\to0$ as $|y|\to\infty$ (Lemma~\ref{lem:decay}), so $W\to0$ at both
ends, and both $R_n''''k_0$ and $R_nk_4$ are integrable on the line. By the fundamental theorem of calculus,
$\int_{\Re z=c}(R_n''''k_0-R_nk_4)\,dz=0$.

\emph{Third equality.} $G_n(c-iy)=\overline{G_n(c+iy)}$, because $R_n$ and $k_4$ commute with complex conjugation.
Hence $\int_{c-i\infty}^{c+i\infty}G_n\,dz=i\int_\RR G_n(c+iy)\,dy=2i\Re\int_0^\infty G_n(c+iy)\,dy$.
\end{proof}

\begin{remark}\label{rem:split}
One can also work with $R_n''''k$ directly: on the upper half of the line split $\pi\cot\pi z=-\pi i+K_+(z)$ with
$K_+(z)=-2\pi i\mathbf q/(1-\mathbf q)$ and integrate the constant part exactly. This gives
$L_n=-R_n'''(c)-\frac1\pi\Im\bigl(\int R_n''''K_+\bigr)$, with a boundary term $-R_n'''(c)$ (an Euler--Maclaurin term)
that has to be bounded separately. The phase is the same: $\sin^5(\pi z)\,e^{2\pi iz}$ has leading exponential
$e^{-3\pi iz}$, which is the origin of the phase $3\pi$ in both routes. Moving the four
derivatives onto the kernel, as above, removes the boundary term, the Fa\`a di Bruno expression for $R_n''''$ and
all asymptotics of the digamma function and its derivatives; only Stirling's formula for $\log\Gamma$ is needed.
\end{remark}

\subsection{Gamma factorisation}

We use the following factors $(c_i,s_i,u_i,\beta_i)$:
\begin{gather*}
  (5,+1,0,0),\quad(5,-1,r,1),\quad(5,+1,1+r,1),\quad(-5,+1,1,1),\\
  \text{and for each block } j:\quad(m_j,+1,\delta_j,0),\quad(-m_j,+1,1-\delta_j,1).
\end{gather*}
Put $\alpha_i(x):=s_ix+u_i$. One checks directly that
\begin{equation}\label{eq:factorid}
  \sum_ic_is_i=0,\qquad\sum_ic_iu_i=-G,\qquad\sum_ic_i\bigl(\beta_i-\tfrac12\bigr)=-S,\qquad\sum_ic_i=10,
\end{equation}
so that $\sum_ic_i\alpha_i(x)=-G$ for all $x$. For Z7 there are $40$ factors, with $\sum_i|c_i|=56$ and
$\sum_i|c_i(\beta_i-\frac12)|=28$.

\begin{lemma}\label{lem:gamma}
For $z\in\CC\setminus\ZZ$,
\[
  R_n(z)=(-1)^{rn}(2z+n)^e\,E_n(z)\,\frac{\sin^5\pi z}{\pi^5},\qquad
  E_n(z):=\prod_i\Gamma(s_iz+nu_i+\beta_i)^{c_i}.
\]
Hence, for $\Im z>0$, $G_n(z)=(-1)^{rn}(2z+n)^eE_n(z)\,T(z)$. For Z7, $(-1)^{rn}=1$ and $e=1$.
\end{lemma}

\begin{proof}
$\prod_{l=1}^{rn}(z-l)=(-1)^{rn}\Gamma(rn+1-z)/\Gamma(1-z)=(-1)^{rn}\Gamma(z)\Gamma(rn+1-z)\sin(\pi z)/\pi$ by the
reflection formula, $\prod_{l=n+1}^{(1+r)n}(z+l)=\Gamma(z+(1+r)n+1)/\Gamma(z+n+1)$, and
$\prod_{k=a_j}^{b_j}(z+k)=\Gamma(z+b_j+1)/\Gamma(z+a_j)$. Raise to the appropriate powers and use
$(-1)^{5rn}=(-1)^{rn}$. The second statement follows from Lemma~\ref{lem:trig}.
\end{proof}

For $\Im z>0$ every argument $s_iz+nu_i+\beta_i$ has nonzero imaginary part, so $E_n$ is holomorphic and zero-free
there, and $G_n$ involves no singular kernel.

\subsection{Stirling's formula in sectors}

Let $\log\Gamma$ denote the branch of the logarithm of $\Gamma$ on $\CC\setminus(-\infty,0]$ that is real on
$(0,\infty)$.

\begin{theorem}[{cf.~\cite[\S5.11(ii)]{DLMF}}]\label{thm:stirling}
Let $0<\varphi\le\pi$ and $\varsigma:=\sin(\varphi/2)$. For $w\ne0$ with $|\arg w|\le\pi-\varphi$,
\[
  \Bigl|\log\Gamma(w)-\bigl((w-\tfrac12)\log w-w+\tfrac12\log2\pi\bigr)\Bigr|\le\frac1{12\varsigma^2|w|}.
\]
\end{theorem}

\begin{proof}
This is the case $n=1$ of the error bound in \cite[\S5.11(ii)]{DLMF}: the remainder of the Stirling series
\cite[(5.11.1)]{DLMF} after its first term is at most the first neglected term $1/(12|w|)$ times
$\sec^2(\frac12\arg w)\le1/\varsigma^2$. We recall the argument, which is the one formalised. With
$Q_1(t):=\frac12\{t\}(1-\{t\})\in[0,\frac18]$ one has
$\log\Gamma(w)=(w-\frac12)\log w-w+\frac12\log2\pi+\int_0^\infty Q_1(t)(t+w)^{-2}\,dt$. If $|\arg w|\le\pi-\varphi$
and $t\ge0$, then $|t+w|\ge\varsigma(t+|w|)$. For $\varrho>0$, the identity
$\int_0^1\frac{u(1-u)/2}{(u+x)^2}du=(x+\frac12)\log(1+\frac1x)-1\le\frac1{12x}-\frac1{12(x+1)}$, summed over the unit
intervals, gives $\int_0^\infty Q_1(t)(t+\varrho)^{-2}dt\le1/(12\varrho)$. Take $\varrho=|w|$.
\end{proof}

\begin{corollary}\label{cor:stirling}
Let $0<\varphi\le\pi$, $\varsigma=\sin(\varphi/2)$, $a>0$, and let $\alpha\in\CC$ with $|\alpha|\ge a$ and
$|\arg\alpha|\le\pi-\varphi$. Let $\beta\in[0,2]$ and $n\ge4/a$. Then
\begin{multline*}
  \Bigl|\log\Gamma(n\alpha+\beta)-\bigl[n\alpha\log n+n(\alpha\log\alpha-\alpha)
   +(\beta-\tfrac12)(\log n+\log\alpha)+\tfrac12\log2\pi\bigr]\Bigr|\\
   \le\frac{9+1/(6\varsigma^2)}{an}.
\end{multline*}
\end{corollary}

\begin{proof}
Put $w=n\alpha+\beta$ and $u=\beta/(n\alpha)$, so $|u|\le2/(an)\le\frac12$ and $|w|\ge an/2$. Adding $\beta\ge0$ moves
the argument towards $0$, so $|\arg w|\le|\arg\alpha|$, and $\log w=\log n+\log\alpha+\log(1+u)$ with principal
logarithms. Insert this into Theorem~\ref{thm:stirling} and use $|\log(1+u)-u|\le|u|^2$ and
$|\log(1+u)|\le\frac32|u|$ for $|u|\le\frac12$.
\end{proof}

\subsection{The phase and the integrand}

On the set of $x\in\CC$ with $\alpha_i(x)\notin(-\infty,0]$ for all $i$ (it contains the open upper half-plane
and the real interval $(\max_j(-\delta_j),r)$) define, with principal logarithms,
\begin{gather*}
  \Psi(x):=\sum_ic_i\,\alpha_i(x)\log\alpha_i(x),\qquad f(x):=\Psi(x)-3\pi ix,\\
  \psi_0(x):=\sum_ic_i(\beta_i-\tfrac12)\log\alpha_i(x),\qquad A(x):=(2x+1)^e\,e^{\psi_0(x)} .
\end{gather*}
By \eqref{eq:factorid}, $\Psi'(x)=\sum_ic_is_i\log\alpha_i(x)$, $\Psi''(x)=\sum_ic_i/\alpha_i(x)$ and
$\Psi'''(x)=-\sum_ic_is_i/\alpha_i(x)^2$. We say that the \emph{sector condition $(a,\eta)$ holds at $x$} if
$|\alpha_i(x)|\ge a$ and $|\arg\alpha_i(x)|\le\pi-\eta$ for every $i$. Put
\[
  \lambda_n:=-Gn\log n+Gn-S\log n+5\log2\pi .
\]

\begin{proposition}\label{prop:integrand}
Let $a>0$ and $0<\eta\le\pi$. There are $n_1$ and a function $B(\varrho)$, depending only on $\cP$, $a$ and $\eta$,
with the following property. Let $n\ge n_1$, and let $x$ satisfy $\Im x>0$, $\Re x\ge0$ and the sector condition
$(a,\eta)$. Then
\[
  G_n(nx)=(-1)^{rn}\,n^e\,e^{\lambda_n+nf(x)}A_n(x)
\]
for a number $A_n(x)$ with
\begin{enumerate}
\item[(U)] $|A_n(x)|\le B(|x|)$;
\item[(M)] $|A_n(x)-A(x)|\le B(|x|)\bigl(n^{-1}+e^{-2\pi n\Im x}\bigr)$.
\end{enumerate}
\end{proposition}

\begin{proof}
Let $\varsigma=\sin(\eta/2)$ and apply Corollary~\ref{cor:stirling} to each factor, with $\alpha=\alpha_i(x)$ and
$\beta=\beta_i$; note $s_i(nx)+nu_i+\beta_i=n\alpha_i(x)+\beta_i$. Summing with weights $c_i$ and using
\eqref{eq:factorid}, the main terms add up to
$-Gn\log n+n\Psi(x)+Gn-S\log n+\psi_0(x)+5\log2\pi=\lambda_n+n\Psi(x)+\psi_0(x)$. Hence
$E_n(nx)=\exp(\lambda_n+n\Psi(x)+\psi_0(x)+e_n)$ with $|e_n|\le\sum_i|c_i|\,(9+1/(6\varsigma^2))/(an)$. By
Lemmas~\ref{lem:gamma} and~\ref{lem:trig}, $(2nx+n)^e=n^e(2x+1)^e$ and
$T(nx)=e^{-3\pi inx}Q(\mathbf q)$ with $Q(\mathbf q)=1+11\mathbf q+11\mathbf q^2+\mathbf q^3$ and
$|\mathbf q|=e^{-2\pi n\Im x}$. So the formula holds with
\[
  A_n(x):=(2x+1)^e\,e^{\psi_0(x)+e_n}\,Q(\mathbf q).
\]
Since $a\le|\alpha_i(x)|\le|x|+|u_i|+2$, $|\psi_0(x)|$ and $|2x+1|$ are bounded in terms of $a$ and $|x|$. With
$|Q|\le24$, $|Q-1|\le23|\mathbf q|$ and $|e^{e_n}-1|\le2|e_n|$ for large $n$, (U) and (M) follow.
\end{proof}

(For Z7 the formalisation proves (U) and (M) in the explicit forms $|A_n|\le72B^2$ and
$|A_n-A|\le100B^2\delta_n$, where $B\ge1$ bounds $|2x+1|$ and $|e^{\psi_0(x)}|$, and $\delta_n\le1$ bounds $|e_n|$
and $|\mathbf q|$.)

\subsection{The certified saddle data}

\begin{theorem}[Certified saddle data, Z7]\label{thm:saddle}
For Z7 put $a:=1/100$, $\eta:=\pi-2.76$, $\rho:=1/20$, $\delta_0:=10^{-7}$, $\etafar:=0.0199414$ and
\[
  c_0:=0.8921643007941758286102076+0.1829120718451090538924822\,i .
\]
There are $x_0\in\CC$, $d\in\CC$ with $|d|=1$, and a polyline $v_0,\dots,v_{23}\in\CC$ with
\begin{gather*}
  v_0=x_c:=0.817411103201,\qquad v_5=x_0-\rho d,\qquad v_6=x_0+\rho d,\\
  v_{23}=\xtop:=2.051154018351+12.001186644670\,i,
\end{gather*}
and $\Re v_k\ge0$ for all $k$, such that the following hold, where $\Omega:=f(x_0)$ and $K:=|\Psi''(x_0)|$.
\begin{enumerate}
\item[(S0)] $|x_0-c_0|\le10^{-12}$, and $x_0$ is the unique zero of $\Psi'-3\pi i$ in this disc. Moreover
  \begin{gather*}
    -10.1167486\le\Re\Omega\le-10.1167485,\\
    -7.704987202\le\Im\Omega\le-7.704987200,\qquad
    \sin(240\,\Im\Omega)<-0.9 .
  \end{gather*}
\item[(S1)] $K\ge28$, $\Psi''(x_0)d^2=-K$, and there is $M\le311$ with $|\Psi'''(x)|\le M$ for $|x-x_0|\le\rho$ and
  $M\rho\le\frac32K$.
\item[(S1$'$)] $\rho<\Im x_0$, so the closed disc $|x-x_0|\le\rho$ lies in $\Im x>0$; and
  $\Im(v_k+t(v_{k+1}-v_k))>0$ for $0\le k<23$ and $0<t\le1$.
\item[(S2)] For $1\le k<23$, $k\ne5$, and every $x\in[v_k,v_{k+1}]$: $\Re f(x)\le\Re\Omega-\etafar$.
\item[(S3)] For every $y\ge0$: $\Re f(\xtop+iy)\le\Re\Omega-\etafar-(12-3\pi)y$.
\item[(S4)] The sector condition $(a,\eta)$ holds on the whole contour: at $x_0+sd$ for $|s|\le\rho$, on the
  segments of (S2), at $\xtop+iy$ for $y\ge0$, and on the segments of (S5).
\item[(S5)] $\delta_0\le x_c<r=\frac{11}{12}$, and for every real $\epsilon$ with $|\epsilon|\le\delta_0$ and every
  $x\ne v_0+\epsilon$ on the segment $[v_0+\epsilon,v_1]$: $\Re f(x)\le\Re\Omega-\etafar$.
\end{enumerate}
\end{theorem}


\begin{proof}[Method of proof]
All statements are verified by exact rational interval arithmetic, with verified enclosures of $\pi$, $\log$,
$\arctan$ and $\arg$ (series with explicit tail bounds after argument reduction) and complex boxes. The verification
is carried out by the Lean kernel; see Section~\ref{sec:lean-cert}.
\begin{itemize}
\item (S0): with $Y\approx1/\Psi''(c_0)$, box evaluations show that $z\mapsto z-Y(\Psi'(z)-3\pi i)$ is a contraction
  of the disc $|z-c_0|\le10^{-12}$ into itself, and the Banach fixed-point theorem gives existence and uniqueness.
  $\Omega$ is enclosed by a box evaluation of $f(c_0)$ (precision $2^{-80}$, $40$ series terms) widened by the
  mean-value bound $|f(x_0)-f(c_0)|\le(10^{-16}+30\cdot10^{-12})\cdot10^{-12}$. For the sine, one checks
  $240\,\Im\Omega=(u-\frac\pi2)-294\cdot2\pi$ with $|u|\le0.38$, so $\sin(240\Im\Omega)=-\cos u\le-(1-u^2/2)<-0.9$.
\item (S1), (S1$'$): $\Psi''(x_0)$ is enclosed in a box; $d:=v/|v|$ with $v:=1-\overline{\Psi''(x_0)}/K$ satisfies
  $\Psi''(x_0)d^2=-K$; and $M$ bounds $\sum_i|c_i|/|\alpha_i|^2$ on the disc, via lower bounds for $|\alpha_i|$
  there. Numerically $K\approx28.31$ and $M\approx310.2$, so $M\rho\approx15.5\le42$; and $\Im x_0\approx0.1829$.
\item (S1$'$), (S2), (S4), (S5): the vertices $v_1,\dots,v_4$ and $v_7,\dots,v_{23}$ are rational points with $12$ decimals, obtained
  from a Runge--Kutta trace of the steepest-descent curves of $\Re f$ from $x_0\pm\rho d$ and a greedy
  simplification ($4$ segments on the lower curve down to $x_c$, $16$ on the upper curve up to $\xtop$). On each
  segment, Taylor bounds on boxes show $\Re f\le T_0:=-10.1367$; since $\Re\Omega\ge-10.1167486$, this gives
  $\Re f\le\Re\Omega-0.0199514$. The connectors $[v_4,v_5]$ and $[v_6,v_7]$ lie in discs of radius $10^{-9}$ around
  $v_4$ and $v_7$, on which $\Re f\le T_0$ is checked directly. Moving the start point by $|\epsilon|\le\delta_0$
  changes $f$ by at most $\delta_0\sup|f'|\le10^{-5}$ on a box containing the widened first segment, which costs
  $10^{-5}$ of the margin; hence $\etafar=0.0199414$. The sector condition is checked on boxes covering all
  segments.
\item (S3): $\partial_y\Re f(\xtop+iy)=3\pi-\Im\Psi'(\xtop+iy)$. Box evaluations give $\Im\Psi'\ge12$ for
  $12.0012\le\Im\le40$, and for $Y:=\Im\xtop+y\ge40$,
  \[
    \Im\Psi'(X+iY)\ge\frac\pi2\sum_ic_i-\sum_i\frac{|c_i||s_iX+u_i|}{Y}=5\pi-\frac{142.86\ldots}{Y}\ge12,
  \]
  with $X:=\Re\xtop$. Since $\Re f(\xtop)\le T_0$, (S3) follows. On the ray $|\alpha_i|\ge12$ and
  $|\arg\alpha_i|\le\frac\pi2+\frac13\le2.76$.\qedhere
\end{itemize}
\end{proof}

These labels are those of the Lean structure \lean{Z7.Analytic.SaddleData}
(\lean{ZetaIrrational\slash{}Z7\slash{}Analytic\slash{}SaddleData.lean}): (S0) is \lean{S0}, \lean{sin\_ne},
\lean{Re$\Omega$\_le}; (S1) is \lean{S1\_dir}, \lean{S1\_M}, \lean{S1\_$\rho$}; (S1$'$) is \lean{S1'\_disc},
\lean{path\_im}; (S2) and (S3) are the fields \lean{S2} and \lean{S3}; (S4) is \lean{S4\_disc} together with the
sector parts of \lean{S2}, \lean{S3} and \lean{S2\_first}; and (S5) is \lean{v0\_im}, \lean{$\delta_0$\_le},
\lean{v0\_lt} and
\lean{S2\_first}. The path fields \lean{v\_disc\_lo}, \lean{v\_disc\_hi}, \lean{re\_nonneg} are the conditions on
the polyline stated before (S0).

At the saddle, $A(x_0)=(2x_0+1)e^{\psi_0(x_0)}\ne0$, since $\Re x_0>0$. The numerical values are
$x_0\approx0.8921643+0.1829121i$, $\Re\Omega\approx-10.11674855$ and $\sin(240\Im\Omega)\approx-0.932453$.

\subsection{Deformation of the contour}

\begin{lemma}\label{lem:deform}
Let $c\in(0,rn+1)\setminus\ZZ$, and let $w_0=c,w_1,\dots,w_{K'}\in\CC$ be such that every point of every segment
$[w_k,w_{k+1}]$ other than $w_0$ has positive imaginary part, and $\Re w_k\ge0$ for all $k$. Then
\[
  i\int_0^\infty G_n(c+iy)\,dy=\sum_{k=0}^{K'-1}\int_{[w_k,w_{k+1}]}G_n(z)\,dz+i\int_0^\infty G_n(w_{K'}+iy)\,dy .
\]
\end{lemma}

\begin{proof}
Let $\delta:=\min_{m\in\ZZ}|c-m|\in(0,\frac12]$; note $\delta\le c$. On $U:=\{\Im z>0\}\cup\{|z-c|<\delta\}$, $G_n$ is
holomorphic: $U$ contains no integer and no pole of $R_n$. $U$ is star-shaped with respect to $c+i\delta/2$, so
$G_n$ has a primitive on $U$. For $Y>\Im w_{K'}$ consider the closed polygon
$c\to c+iY\to\Re w_{K'}+iY\to w_{K'}\to w_{K'-1}\to\dots\to w_0=c$, which lies in $U$; the integral of $G_n$ around it
is $0$. On the horizontal side, which lies in $\Re z\ge0$, $k_4$ is bounded and $|R_n|\le C(1+|z|^2)^{-1}$ by
Lemma~\ref{lem:decay}, so its contribution tends to $0$ as $Y\to\infty$. The vertical integrals converge absolutely.
\end{proof}

\subsection{Laplace's method}

\begin{lemma}\label{lem:laplace}
Let $K_1,\rho_1>0$, $M_1\ge0$ with $M_1\rho_1\le\frac32K_1$. Let $\phi,g:[-\rho_1,\rho_1]\to\CC$ be continuous,
with $|\phi(s)+K_1s^2/2|\le M_1|s|^3/6$, $|g|\le A_0$ and $|g(s)-g(0)|\le A_1|s|$. Then for $n>0$
\[
  \Bigl|\int_{-\rho_1}^{\rho_1}e^{n\phi(s)}g(s)\,ds-g(0)\sqrt{\frac{2\pi}{nK_1}}\Bigr|
  \le\frac1n\Bigl(\frac{8A_0M_1}{3K_1^2}+\frac{2A_1}{K_1}\Bigr)+A_0e^{-nK_1\rho_1^2/4}\sqrt{\frac{4\pi}{nK_1}} .
\]
Moreover $\Re\phi(s)\le0$.
\end{lemma}

(The formalisation proves the same estimate with the larger constants $(16\sqrt{2\pi}/3)\,A_0M_1/K_1^2$ and
$2\sqrt{2\pi}\,A_1/K_1$; only the order $O(1/n)$ is used.)

\begin{proof}
Write $\phi(s)=-K_1s^2/2+\tau(s)$. Then $|\tau(s)|\le M_1|s|^3/6\le M_1\rho_1s^2/6\le K_1s^2/4$, so
$\Re\phi\le-K_1s^2/4\le0$ and
$|e^{n\phi}-e^{-nK_1s^2/2}|\le e^{-nK_1s^2/2}n|\tau|e^{n|\tau|}\le\frac{nM_1|s|^3}6e^{-nK_1s^2/4}$. Split
$\int e^{n\phi}g=g(0)\int e^{-nK_1s^2/2}+\int e^{-nK_1s^2/2}(g(s)-g(0))+\int(e^{n\phi}-e^{-nK_1s^2/2})g$. The second
integral is at most $A_1\int_\RR|s|e^{-nK_1s^2/2}ds=2A_1/(nK_1)$, and the third at most
$\frac{nM_1A_0}6\int_\RR|s|^3e^{-nK_1s^2/4}ds=\frac{8A_0M_1}{3nK_1^2}$. Finally
$\int_{|s|>\rho_1}e^{-nK_1s^2/2}ds\le e^{-nK_1\rho_1^2/4}\int_\RR e^{-nK_1s^2/4}ds$.
\end{proof}

\subsection{The asymptotic formula, nonvanishing and size}

\begin{proposition}\label{prop:asymp}
For Z7 let $\theta:=\Im\Omega$ and $C_0:=d\,A(x_0)\sqrt{2\pi/K}$, so $C_0\ne0$; write $C_0=|C_0|e^{i\phi_0}$. Then,
as $n\to\infty$ in $240\NN$,
\[
  L_n=-\frac1\pi\,n^{3/2}e^{\lambda_n+n\Re\Omega}\bigl(|C_0|\sin(n\theta+\phi_0)+o(1)\bigr).
\]
\end{proposition}

\begin{proof}
Let $c_n:=\lfloor nx_c\rfloor+\frac12$. Since $0<x_c<r$, we have $0<c_n<rn+1$ and $c_n\notin\ZZ$, and
$\chi_n:=c_n/n-x_c$ satisfies $|\chi_n|\le\frac1{2n}\le\delta_0$ for large $n$. Let
$J_n:=i\int_0^\infty G_n(c_n+iy)\,dy$. By Proposition~\ref{prop:contour}, $L_n=-\frac1\pi\Im J_n$. Apply
Lemma~\ref{lem:deform} with $w_0=c_n$ and $w_k=nv_k$ $(1\le k\le23)$; the hypotheses hold by (S1$'$) and $\Re v_k\ge0$, since
$\Im(c_n+t(nv_1-c_n))=tn\Im v_1>0$ for $t>0$. Thus $J_n$ is the sum of the $23$ segment integrals and the integral
over the vertical ray from $n\xtop$.

\emph{The Taylor segment.} Parametrising $[nv_5,nv_6]$ by $z=n(x_0+sd)$, $|s|\le\rho$, and using
Proposition~\ref{prop:integrand} (applicable by (S1$'$), (S4) and $\Re(x_0+sd)\ge0$, which holds since
$\Re v_5,\Re v_6\ge0$),
\[
  \int_{[nv_5,nv_6]}G_n(z)\,dz=n^2d\,e^{\lambda_n+n\Omega}\mathcal D_n,\qquad
  \mathcal D_n:=\int_{-\rho}^{\rho}e^{n\phi(s)}A_n(x_0+sd)\,ds,
\]
where $\phi(s):=f(x_0+sd)-\Omega$. Since $f'(x_0)=\Psi'(x_0)-3\pi i=0$, $f''(x_0)d^2=\Psi''(x_0)d^2=-K$ and
$|f'''|=|\Psi'''|\le M$ on the disc, Taylor's formula gives $|\phi(s)+Ks^2/2|\le M|s|^3/6$; by
Lemma~\ref{lem:laplace}, $\Re\phi\le0$. By (M) with $\Im(x_0+sd)\ge\Im x_0-\rho>0$,
$|\mathcal D_n-\int_{-\rho}^{\rho}e^{n\phi(s)}A(x_0+sd)\,ds|=O(n^{-1})$. The function $A$ is holomorphic on $\Im x>0$,
hence bounded and Lipschitz on the disc, and Lemma~\ref{lem:laplace} gives
$\mathcal D_n=A(x_0)\sqrt{2\pi/(nK)}+O(n^{-1})$.

\emph{The far part.} Let $\mathcal E_n$ be the sum of the other $22$ segment integrals and the ray integral. On each
of these segments, (S2), (S4) and (U) give $|G_n(z)|\le Cn\,e^{\lambda_n+n(\Re\Omega-\etafar)}$, and the segments
have length $O(n)$. (For the first segment, $z=n((x_c+\chi_n)+t(v_1-x_c-\chi_n))$ with $|\chi_n|\le\delta_0$, and (S5)
applies.) On the ray, $z=n(\xtop+iu)$ with $u\ge0$, and (S3), (S4), (U) give
$|G_n(z)|\le n\,e^{\lambda_n+n(\Re\Omega-\etafar)}e^{-n(12-3\pi)u}B(|\xtop|+u)$; since $12-3\pi>2.5$, its integral over
$dz=in\,du$ is $O(n^2e^{\lambda_n+n(\Re\Omega-\etafar)})$ for large $n$. Hence
$|\mathcal E_n|=O(n^2e^{\lambda_n+n\Re\Omega-n\etafar})$.

\emph{Conclusion.} Combining,
\[
\begin{split}
  n^{-3/2}e^{-\lambda_n}J_n&=e^{n\Omega}\bigl(C_0+O(n^{-1/2})\bigr)+O\bigl(n^{1/2}e^{-n\etafar}\bigr)e^{n\Re\Omega}\\
  &=C_0e^{n\Omega}+o\bigl(e^{n\Re\Omega}\bigr),
\end{split}
\]
and $\Im(C_0e^{n\Omega})=|C_0|e^{n\Re\Omega}\sin(n\theta+\phi_0)$.
\end{proof}

\begin{lemma}\label{lem:sines}
For real $u,v$: $|\sin u|+|\sin(u+v)|\ge|\sin v|$.
\end{lemma}

\begin{proof}
$\sin v=\sin(u+v)\cos u-\cos(u+v)\sin u$.
\end{proof}

\begin{corollary}\label{cor:nonvan}
For Z7 and all sufficiently large $n\in240\NN$, $L_n\ne0$ or $L_{n+240}\ne0$. In particular $L_n\ne0$ for
infinitely many $n\in240\NN$.
\end{corollary}

\begin{proof}
By Lemma~\ref{lem:sines} with $v=240\theta$ and (S0),
$\max(|\sin(n\theta+\phi_0)|,|\sin((n+240)\theta+\phi_0)|)\ge\frac12|\sin(240\theta)|>0.45$ for every $n$. By
Proposition~\ref{prop:asymp}, once the $o(1)$ terms are below $0.45|C_0|$, the corresponding $L$ is nonzero.
\end{proof}

\begin{proposition}[Size]\label{prop:size}
For Z7 put $\Vcert:=-10.1167$. There are $\eps''_n\to0$ such that, for all large $n\in240\NN$ with $L_n\ne0$,
\[
  \log|L_n|\le-G\log n!+n\bigl(\Vcert+\eps''_n\bigr).
\]
\end{proposition}

\begin{proof}
By Proposition~\ref{prop:contour} and the proof of Proposition~\ref{prop:asymp},
$|L_n|\le|J_n|/\pi\le Cn^{3/2}e^{\lambda_n+n\Re\Omega}$ for large $n$. By Stirling's upper bound
$\log n!\le n\log n-n+\frac12\log n+1$,
$\lambda_n\le-G\log n!+G(\frac12\log n+1)-18\log n+5\log2\pi$. Finally $\Re\Omega\le\Vcert$ by (S0).
\end{proof}

\begin{remark}\label{rem:realaxis}
In the $\beta$-analogue \cite{BetaPaper} the size of $L_n$ was bounded by an elementary estimate on the real axis.
Here that route is too weak: on the real axis the integrand of $L_n$ only gives the exponent
$V_{\mathrm{up}}\approx-8.92$ for Z7 [numerical], against $\Acert\approx9.92$, so that $V_{\mathrm{up}}+\Acert\approx+1.0$.
The saddle point gains about $1.19$ per unit of $n$, through the cancellation in the sum $\sum_tR_n''''(t)$, and this
gain is essential for Z7. (For Z6 the real-axis exponent would leave a margin of only about $0.07$; see
Table~\ref{tab:certs}.)
\end{remark}

% ===========================================================================
\section{Proof of Theorem~\ref{thm:main}}\label{sec:proof}
% ===========================================================================

Put
\begin{equation}\label{eq:Q}
\begin{split}
  Q:=\Vcert+\Acert&=-10.1167+9.9185802528\ldots\\
  &=-0.1981197\ldots<0 .
\end{split}
\end{equation}
(With the certified upper end $-10.1167485$ of $\Re\Omega$ in place of $\Vcert$ one gets $Q\le-0.1981682$.) Let
$\ell_n:=D_nL_n$, $n\in240\NN$. By Corollary~\ref{cor:intform}, $\ell_n=a_n+\sum_{j=3}^{10}c_{n,j}\zeta(2j+1)$ with
$a_n,c_{n,j}\in\ZZ$. If $L_n\ne0$, Corollary~\ref{cor:logD} and Proposition~\ref{prop:size} give, for large $n$,
\[
  \log|\ell_n|=\log D_n+\log|L_n|\le n\bigl(Q+\eps'_n+\eps''_n\bigr),
\]
the terms $\pm G\log n!$ cancelling. Since $Q<0$ and $\eps'_n,\eps''_n\to0$, it follows that $|\ell_n|\le e^{nQ/2}$ for
all large $n$ (trivially when $L_n=0$), so $\ell_n\to0$. On the other hand $D_n>0$, so by Corollary~\ref{cor:nonvan},
$\ell_n\ne0$ for infinitely many $n$.

Suppose now that $\zeta(7),\zeta(9),\dots,\zeta(21)$ are all rational, and let $\Delta\ge1$ be a common denominator.
Then $\Delta\ell_n\in\ZZ$ for every $n$. For the infinitely many $n$ with $\ell_n\ne0$ this gives $|\Delta\ell_n|\ge1$,
which contradicts $\ell_n\to0$. Hence at least one of $\zeta(7),\dots,\zeta(21)$ is irrational. \qed

% ===========================================================================
\section{The general theorem and other instances}\label{sec:general}
% ===========================================================================

The proofs of Sections~\ref{sec:construction}--\ref{sec:proof} use the specific parameters of Z7 only through the
following hypotheses.

\begin{theorem}\label{thm:general}
Let $\cP$ be a parameter set (Definition~\ref{def:param}), and suppose:
\begin{enumerate}
\item[(H1)] $G>0$;
\item[(H2)] $\Acert$ is a rational number with
  $\sum_{\ell}e^*(\xi_\ell)(\frac1{x_{\ell-1}}-\frac1{x_\ell})+\frac{S+15}{X}+\frac{G}{2(X-1)}\le\Acert$, for the event
  set of Section~\ref{sec:A} with $X=400$;
\item[(H3)] there are saddle data $(x_0,d,K,M,\rho,v_0,\dots,v_{K_{\mathrm{top}}})$ and constants
  $a,\eta,\etafar,\delta_0>0$ satisfying the analogues of (S0)--(S5) of Theorem~\ref{thm:saddle} (without the
  numerical values), with $\Re v_k\ge0$, a disc segment $[v_{K_d},v_{K_d+1}]=[x_0-\rho d,x_0+\rho d]$, and some
  positive decay rate on the top ray in (S3);
\item[(H4)] $\sin(q\Im\Omega)\ne0$;
\item[(H5)] $Q:=\Re\Omega+\Acert<0$;
\item[(H6)] $qG\in2\ZZ$, and $e\equiv1+S\pmod2$.
\end{enumerate}
Then at least one of $\zeta(7),\zeta(9),\dots,\zeta(N_\cP)$ is irrational.
\end{theorem}

\begin{proof}
Sections~\ref{sec:algebra}--\ref{sec:A} are written for a general parameter set; (H1) and (H6) are what they use
(Lemma~\ref{lem:sym} and the remark after it), and (H2) replaces Proposition~\ref{prop:certA}, the tail constant
being $S+15$ by Lemma~\ref{lem:tail}. Section~\ref{sec:nonvan} uses only (H3), with $\lambda_n$ and the amplitude
$A(x)=(2x+1)^ee^{\psi_0(x)}$ as defined there, and the factor $(-1)^{rn}$ of Lemma~\ref{lem:gamma} does not affect
the argument, which uses only $|\sin(n\theta+\phi_0)|$; (H4) replaces the sine bound of (S0). The asymptotic
formula becomes $L_n=\mp\frac1\pi n^{e+1/2}e^{\lambda_n+n\Re\Omega}(|C_0|\sin(n\theta+\phi_0)+o(1))$. Section~\ref{sec:proof}
then needs (H5).
\end{proof}

Theorem~\ref{thm:main} is Theorem~\ref{thm:general} for Z7, and this is the case that is formalised. Seven
instances have been certified; Tables~\ref{tab:instances} and~\ref{tab:certs} list them, with all bounds rounded in
the safe direction. Z1 has a single block and reaches Zudilin's window $\zeta(7),\dots,\zeta(37)$; Z4--Z7 have
near-continuous profiles.

\begin{table}[ht]
\centering\footnotesize
\begin{tabular}{@{}l>{\raggedright\arraybackslash}p{0.07\textwidth}l>{\raggedright\arraybackslash}p{0.43\textwidth}llll@{}}
\toprule
 & $r$ & $q$ & blocks $\delta{:}\,m$ & $S$ & $N_\cP$ & $G$; $qG$ & $e$\\
\midrule
Z1 & $5/4$ & $4$ & $0{:}\,33$ & $33$ & $37$ & $41/2$; $82$ & $0$\\
Z2 & $1$ & $8$ & $0{:}\,21$, $\frac18{:}\,4$, $\frac14{:}\,4$ & $29$ & $33$ & $16$; $128$ & $0$\\
Z3 & $1$ & $8$ & $0{:}\,19$, $\frac18{:}\,4$, $\frac14{:}\,4$ & $27$ & $31$ & $14$; $112$ & $0$\\
Z4 & $1$ & $80$ & $0{:}\,2$, $\frac1{80}{:}\,1$, $\frac1{40}{:}\,2$, $\frac3{80}{:}\,1$, $\frac1{20}{:}\,1$,
  $\frac1{16}{:}\,1$, $\frac3{40}{:}\,1$, $\frac7{80}{:}\,1$, $\frac1{10}{:}\,1$, $\frac18{:}\,2$, $\frac{11}{80}{:}\,1$,
  $\frac{13}{80}{:}\,2$, $\frac7{40}{:}\,1$, $\frac3{16}{:}\,1$, $\frac15{:}\,1$, $\frac{17}{80}{:}\,1$,
  $\frac9{40}{:}\,1$, $\frac14{:}\,2$, $\frac{23}{80}{:}\,2$ & $25$ & $29$ & $339/40$; $678$ & $0$\\
Z5 & $1$ & $80$ & $0{:}\,2$; $\frac k{80}{:}\,1$, $k=1,\dots,19$ & $21$ & $25$ & $25/4$; $500$ & $0$\\
Z6 & $19/24$ & $240$ & $\frac k{80}{:}\,1$, $k=1,\dots,19$ & $19$ & $23$ & $19/3$; $1520$ & $0$\\
Z7 & $11/12$ & $240$ & $\frac k{80}{:}\,1$, $k=1,\dots,18$ & $18$ & $21$ & $547/120$; $1094$ & $1$\\
\bottomrule
\end{tabular}
\caption{The certified instances: parameters. Each gives ``one of $\zeta(7),\zeta(9),\dots,\zeta(N_\cP)$''.}\label{tab:instances}
\end{table}

\begin{table}[ht]
\centering\footnotesize\setlength{\tabcolsep}{2pt}
\begin{tabular}{@{}lllllll@{}}
\toprule
 & $\Acert\le$ & $\Re\Omega\le$ & $Q\le$ & $\sin(q\Im\Omega)$ & $x_0$ & $V_{\mathrm{up}}$\\
\midrule
Z1 & $35.7955767$ & $-36.7604812$ & $-0.9649$ & $-0.709656$ & $1.2579692+0.0404800i$ & $-36.597$\\
Z2 & $25.3852394$ & $-26.4996942$ & $-1.1144$ & $0.997966$ & $1.0073349+0.0420146i$ & $-26.327$\\
Z3 & $23.3827332$ & $-23.7154311$ & $-0.3326$ & $0.825398$ & $1.0077298+0.0558821i$ & $-23.476$\\
Z4 & $11.4875113$ & $-16.0979480$ & $-4.6104$ & $-0.968332$ & $0.9994996+0.1230432i$ & $-15.461$\\
Z5 & $10.8538930$ & $-12.9889571$ & $-2.1350$ & $0.283577$ & $0.9869778+0.1636877i$ & $-12.027$\\
Z6 & $10.7587916$ & $-11.4775814$ & $-0.7187$ & $-0.272087$ & $0.7869951+0.1190009i$ & $-10.826$\\
Z7 & $9.9185803$ & $-10.1167485$ & $-0.1981$ & $-0.932453$ & $0.8921643+0.1829121i$ & $-8.923$\\
\bottomrule
\end{tabular}

\medskip
\begin{tabular}{@{}lllllll@{}}
\toprule
 & (S1) $\kappa$, $M$, $\rho$ & (S1$'$) & (S2) $\etafar$ & (S3) $\Im\Psi'$ & (S4) $a$, $\eta$ & $\eta_{\mathrm{cr}}$; $x_c$\\
\midrule
Z1 & $118.54$, $5527$, $0.01031$ & $0.0302$ & $5.65\cdot10^{-3}$ & $15.47$ & $1.92\cdot10^{-2}$, $0.672$ & $0.2211$; $1.2290$\\
Z2 & $114.57$, $5128$, $0.01066$ & $0.0314$ & $5.82\cdot10^{-3}$ & $15.53$ & $1.99\cdot10^{-2}$, $0.718$ & $0.2256$; $0.9783$\\
Z3 & $86.68$, $2904$, $0.01410$ & $0.0418$ & $7.68\cdot10^{-3}$ & $15.55$ & $2.66\cdot10^{-2}$, $0.788$ & $0.2923$; $0.9712$\\
Z4 & $40.68$, $622.6$, $0.03076$ & $0.0923$ & $1.69\cdot10^{-2}$ & $15.59$ & $6.06\cdot10^{-2}$, $1.029$ & $0.5898$; $0.9357$\\
Z5 & $31.19$, $364.8$, $0.04105$ & $0.1226$ & $2.30\cdot10^{-2}$ & $15.62$ & $8.29\cdot10^{-2}$, $1.138$ & $0.7565$; $0.9128$\\
Z6 & $42.43$, $670.7$, $0.02977$ & $0.0892$ & $1.65\cdot10^{-2}$ & $15.63$ & $5.93\cdot10^{-2}$, $1.077$ & $0.5602$; $0.7290$\\
Z7 & $28.31$, $301.4$, $0.04614$ & $0.1368$ & $2.63\cdot10^{-2}$ & $15.64$ & $9.48\cdot10^{-2}$, $1.210$ & $0.8265$; $0.8174$\\
\bottomrule
\end{tabular}
\caption{The certified instances: certificates of the Python and Arb computations. $\Acert$ is the exact rational
of (H2), rounded up; $\Re\Omega$ is the upper end of the Arb enclosure (radius $\le2\cdot10^{-7}$), rounded up;
$Q=\Re\Omega+\Acert$, rounded up. Second part: the Taylor-disc data $\kappa=-\Re(\Psi''(x_0)d^2)$, $M$, $\rho$; the margin
$\Im x_0-\rho$ of (S1$'$); the far margin; the tail value of $\Im\Psi'$ on the top ray (it must exceed $3\pi$); the
sector constants; and the crossing gap $\eta_{\mathrm{cr}}=\Re\Omega-\Psi(x_c)$ with the start point $x_c$ (a cross-check,
not needed by the proof). $V_{\mathrm{up}}$ is
the real-axis exponent of the numerical model [numerical].}\label{tab:certs}
\end{table}

The exact values of $\Acert$ are
\begin{align*}
\text{Z1}&:\ \tfrac{28289277235925350790655713667586667}{790300921079786767495600904601600},\\
\text{Z2}&:\ \tfrac{20061978064991101891984267756512101}{790300921079786767495600904601600},\\
\text{Z3}&:\ \tfrac{1319956822769190839182373382010337}{56450065791413340535400064614400},\\
\text{Z4}&:\ \tfrac{75654922793088526827628553510069}{6585841008998223062463340871680},\\
\text{Z5}&:\ \tfrac{4288920793391182192917685241907847}{395150460539893383747800452300800},\\
\text{Z6}&:\ \tfrac{671264435316850956055216565463991}{62392177979983165854915860889600},
\end{align*}
and, for Z7, \eqref{eq:Acert}.

\emph{Status of the instances.} For Z7 everything is kernel-checked (Section~\ref{sec:lean}); the Lean saddle
certificate of Theorem~\ref{thm:saddle} uses its own contour (with $\rho=1/20$ and $23$ segments) and agrees with the
Arb value of $\Re\Omega$ in Table~\ref{tab:certs}, and the Lean value of $\Acert$ is the same exact rational. For
Z1--Z6 the certificates come from two programs that are not part of the repository: an exact-rational Python
implementation (python-flint) of the sum in (H2), and a ball-arithmetic computation with Arb of the saddle data for
(H3), (H4). The latter checks a variant of (S0)--(S5): the Taylor radius is chosen with $\rho\le\kappa/M$, which is
stronger than $M\rho\le\frac32\kappa$, and (S2) is checked on the whole polyline starting at $x_c$, together with the
crossing gap $\Psi(x_c)\le\Re\Omega-\eta_{\mathrm{cr}}$; the perturbation of the start point in (S5) is not part of
that certificate. That perturbation changes $\Re f$ on the first segment by $O(1/n)$, because the sector condition
holds on a neighbourhood of that segment, so it costs nothing in the limit [proof sketch]. So Z1--Z6 rest on these
certificates and are not formalised. The corresponding windows $\zeta(7),\dots,\zeta(N)$ for $N=23,\dots,37$ are weaker than
Theorem~\ref{thm:main} in any case; their only interest is the larger margins.

As an independent check of the arithmetic part, for every profile type of Table~\ref{tab:instances} and for Z7
itself at $n=240$, the exact content of the vector $(\rho_n,\tau_i)$ was computed with python-flint and compared with
the exponents $E_p$ for $\sqrt\cN<p\le n$: the bound was never violated, it was attained at most primes (at $30$ of
the $39$ primes $41\le p\le233$ for Z7 at $n=240$), and $E_p=0$ for $p>n$ [numerical]. At the same time the window
was confirmed to be exactly $\zeta(7),\dots,\zeta(N_\cP)$, and $L_n$ was compared with the direct sum
$\sum_tR_n''''(t)$ to $30$ digits.

% ===========================================================================
\section{The formalisation}\label{sec:lean}
% ===========================================================================

\subsection{What is formalised}
The repository \url{https://github.com/telaviv/zeta-irrationality} \cite{ZetaRepo}, at commit \lean{1b48910},
contains a proof of
Theorem~\ref{thm:main} in Lean~4 \cite{LeanFour} (version v4.35.0) with Mathlib \cite{Mathlib} (version v4.35.0
\cite{Mathlib4}). All declaration names and file paths in this section refer to that commit; paths are relative
to the repository root. The formalisation covers every step of the proof of
Theorem~\ref{thm:main}, including the parts that a written proof would cite: the prime number theorem, Stirling's
formula with explicit remainder in sectors, the contour-integral representation, the deformation of contours, and all
numerical certificates. Much of the instance-independent infrastructure (prime number theorem, Stirling, interval
arithmetic, contour and Laplace lemmas, partial fractions, valuations) is adapted from the formalisation of the
$\beta$-analogue \cite{BetaPaper}.

\subsection{The statement}
The definition that the reader has to check is, in \lean{ZetaIrrational\slash{}Statement.lean} (namespace
\lean{ZetaIrr}):
\begin{quote}\ttfamily\footnotesize
def OneOfZeta7To21 : Prop :=\\
\hspace*{2em}$\exists$ j $\in$ (\char123 7,~9,~11,~13,~15,~17,~19,~21\char125~:~Finset~$\mathbb{N}$),\\
\hspace*{4em}Irrational (riemannZeta j).re\\[2pt]
theorem ZetaIrr.Z7.oneOfZeta7To21 : OneOfZeta7To21
\end{quote}
The theorem is in \lean{ZetaIrrational\slash{}Z7\slash{}MainUnconditional.lean}. Here \lean{riemannZeta : $\mathbb C\to\mathbb C$} is
Mathlib's Riemann zeta function, the natural number \lean{j} is cast to $\CC$, and \lean{Irrational x} is Mathlib's
predicate $x\notin\QQ$. Nothing defined in the project occurs in the statement. Mathlib's
\lean{zeta\_nat\_eq\_tsum\_of\_gt\_one} identifies \lean{riemannZeta j} with $\sum_{m\ge1}m^{-j}$ for $j\ge2$; the
file \lean{ZetaIrrational\slash{}Defs.lean} uses it to pass to the real series used internally
(\lean{riemannZeta\_re\_eq\_zetaR}). The statement assumes nothing about $\zeta(3)$ or $\zeta(5)$.

The command \lean{lake env lean Axioms.lean} prints
\begin{quote}\ttfamily\small
'ZetaIrr.Z7.oneOfZeta7To21' depends on axioms:\\
\hspace*{2em}[propext, Classical.choice, Quot.sound]
\end{quote}
and the same axiom list for each of the main ingredients listed in \lean{Axioms.lean}. These are Lean's three
standard axioms, which Mathlib also uses. The development contains no \lean{sorry}, no added axiom, no
\lean{native\_decide} and no \lean{implemented\_by}, so no step relies on compiled code.

\subsection{Correspondence with the paper}
Table~\ref{tab:lean} lists the main formal counterparts. The theorem is \lean{oneOfZeta7To21\_of\_certUnits}
applied to the kernel-checked sum \lean{CertA.bodyUnits\_400}. It goes through \lean{oneOfZeta7To21\_of\_parts},
whose inputs are the certificate inequality, Proposition~\ref{prop:size} and Corollary~\ref{cor:nonvan}, and through
the criterion \lean{exists\_irrational\_of\_small\_forms}, which formalises Section~\ref{sec:proof}.

\begin{table}[ht]
\centering\small
\begin{tabular}{@{}>{\raggedright\arraybackslash}p{0.27\textwidth}>{\raggedright\arraybackslash}p{0.69\textwidth}@{}}
\toprule
paper & Lean declaration (namespace \lean{ZetaIrr}) \\
\midrule
Lemma~\ref{lem:pf}, Prop.~\ref{prop:expansion} & \lean{Z7.R\_eq\_pfR}, \lean{PF.pf\_eval}, \lean{PFData.tsum\_R4}\\
Lemmas~\ref{lem:tau-even}, \ref{lem:tau1}, Cor.~\ref{cor:form} & \lean{Z7.pf\_symmetric}, \lean{PFData.tau\_eq\_zero\_of\_even}, \lean{PF.pf\_top\_coeff}, \lean{Z7.pf\_tau\_one}, \lean{PFData.tsum\_R4\_odd}, \lean{Z7.expansion\_Z7}\\
Lemmas~\ref{lem:psge}, \ref{lem:A}, Cor.~\ref{cor:A} & \lean{Arith.psge\_mul}, \lean{Arith.psge\_inv}, \lean{Z7.Arith.vge\_a}, \lean{Z7.Arith.a\_int\_of\_nodvd}, \lean{Z7.Arith.a\_int}, \lean{Z7.Arith.a\_bound\_sharp}, \lean{Z7.Arith.a\_bound\_small}\\
Lemma~\ref{lem:B} & \lean{Arith.vge\_Hq}, \lean{Arith.vge\_Hq\_of\_lt}\\
Def.~\ref{def:Ep}, Prop.~\ref{prop:integral} & \lean{Z7.Arith.Ep}, \lean{Z7.Arith.D}, \lean{Z7.Arith.int\_rho}, \lean{Z7.Arith.int\_tau}\\
Props.~\ref{prop:smallprimes}, \ref{prop:decomp} & \lean{Z7.Arith.Ep\_sub\_le\_small}, \lean{Z7.Arith.small\_primes\_sum\_le}, \lean{Z7.Arith.log\_D\_sub\_le}\\
Prop.~\ref{prop:pointwise}, Lemma~\ref{lem:tail} & \lean{Z7.CertA.pointwise\_le\_eStar\_prime}, \lean{Z7.Esharp\_sub\_le\_eStar}, \lean{Z7.CertA.eStar\_le\_tail}\\
Lemma~\ref{lem:step} & \lean{Z7.CertA.eStar\_eq\_mid}, \lean{Z7.CertA.eStar\_isGreatest}\\
Prop.~\ref{prop:certA} & \lean{Z7.CertA.bodyUnits\_400}, \lean{Z7.CertA.certFact\_of\_units}\\
Lemma~\ref{lem:pnt-tail}, Prop.~\ref{prop:lemma5} & \lean{PNT.eventually\_primeSum\_lemma5\_le}, \lean{Z7.sharp\_of\_pointwise}, \lean{Z7.arith\_Z7}\\
prime number theorem & \lean{PNT.theta\_div\_tendsto}\\
Lemma~\ref{lem:trig} & \lean{Analytic.hasDerivAt\_kk}, \lean{Analytic.sin\_pow\_five\_mul\_kk4}, \lean{Analytic.T\_eq\_exp}\\
Prop.~\ref{prop:contour} & \lean{Analytic.hasSum\_cot\_eq\_line\_integral}, \lean{Analytic.integral\_line\_by\_parts4}, \lean{Z7.Analytic.L\_eq\_half\_line}\\
Lemma~\ref{lem:gamma}, Prop.~\ref{prop:integrand} & \lean{Z7.Analytic.RC\_eq\_En}, \lean{Z7.Analytic.Gfn\_eq\_PT}, \lean{Z7.Analytic.Pfun\_asymp}, \lean{Z7.Analytic.Gfn\_path}, \lean{Z7.Analytic.norm\_Gfn\_le}, \lean{Z7.Analytic.Gfn\_disc\_approx}\\
Thm.~\ref{thm:stirling}, Cor.~\ref{cor:stirling} & \lean{Stirling.norm\_J\_le}, \lean{Stirling.stirling\_scaled}\\
Thm.~\ref{thm:saddle} & \lean{Z7.Interval.Z7\_saddleData} (a term of the structure \lean{Z7.Analytic.SaddleData})\\
Lemmas~\ref{lem:deform}, \ref{lem:laplace} & \lean{Z7.Analytic.contour\_eq\_polyline\_Gfn}, \lean{Analytic.laplace\_bound}\\
Prop.~\ref{prop:asymp}, Cor.~\ref{cor:nonvan} & \lean{Z7.Analytic.nonvan\_of\_saddleData}, \lean{Z7.nonvan\_Z7}\\
Prop.~\ref{prop:size} & \lean{Z7.Analytic.size\_of\_saddleData}, \lean{Z7.size\_Z7}\\
Section~\ref{sec:proof} & \lean{exists\_irrational\_of\_small\_forms}, \lean{oneOfZeta7To21\_of\_forms}, \lean{Z7.oneOfZeta7To21\_of\_hyp}, \lean{Z7.oneOfZeta7To21\_of\_parts}, \lean{Z7.oneOfZeta7To21\_of\_certUnits}, \lean{Z7.oneOfZeta7To21}\\
\bottomrule
\end{tabular}
\caption{Formal counterparts of the main steps.}\label{tab:lean}
\end{table}

\subsection{Notes on the formal proof}\label{sec:lean-diff}
The formal proof is the proof of this paper. The following points are worth noting when reading the two side by
side.
\begin{itemize}
\item \emph{The analytic part} uses the four integrations by parts of Proposition~\ref{prop:contour}, not the kernel
  split of Remark~\ref{rem:split}. There is no boundary term, and only Stirling's formula for $\log\Gamma$ is needed,
  for shifts $\beta\in[0,2]$ (the factor list contains $\Gamma(z)$ and $\Gamma(z+\delta_jn)$, with $\beta=0$). The
  start of the contour near the real axis is handled by (S5); no separate condition at the crossing point $x_c$ is
  needed, because the integrand $G_n$ has no singular kernel.
\item \emph{The size of $L_n$} comes from the saddle point (Proposition~\ref{prop:size}), with
  $\Vcert=-10.1167$; see Remark~\ref{rem:realaxis}.
\item \emph{Corollary~\ref{cor:A}(b$'$)} is proved without any condition on $p$, and only its integrality form is
  used.
\item \emph{The pointwise bound} (Proposition~\ref{prop:pointwise}) holds for every prime $p\le n$.
\item \emph{Small primes.} Definition~\ref{def:Ep} uses $(S+4)\lambda_p$ (for Z7, $22\lambda_p$) and drops the
  $e$-factor from $\nu_p(k)$; Proposition~\ref{prop:smallprimes} has the explicit constant $182$.
\item \emph{The tail.} Lemma~\ref{lem:tail} gives $e^*(x)\le S+14+G\{x\}$; the certificate uses the slightly larger
  constant $S+15=33$, which costs $1/400$ in $\Acert$.
\item \emph{The event set} for Z7 has $30{,}138$ points, so the sum of Proposition~\ref{prop:certA} has $30{,}137$
  terms.
\item \emph{The saddle certificate} uses its own steepest-descent contour (Theorem~\ref{thm:saddle}: $\rho=1/20$,
  $23$ segments, $\etafar=0.0199414$, $a=1/100$, top ray from $\xtop\approx2.0512+12.0012i$), generated independently
  of the Arb computation of Table~\ref{tab:certs}; the two enclosures of $\Re\Omega$ agree.
\end{itemize}
Further, as for the $\beta$-analogue, a few standard tools are organised differently because Mathlib (at the
version used) lacks them. Mathlib has no residue theorem, so the first equality of Proposition~\ref{prop:contour}
is proved by subtracting model principal parts at the integers $c<t\le N$, removing the resulting removable
singularities, applying Cauchy's theorem for rectangles and evaluating the model integrals explicitly. The deformation
(Lemma~\ref{lem:deform}) uses primitives of holomorphic functions on unions of discs. Stirling's formula is proved
from Euler's limit formula for $\Gamma$, an exact Euler--Maclaurin identity and Mathlib's Stirling formula for
$n!$. The prime number theorem in the form $\theta(y)/y\to1$ is ported from the PrimeNumberTheoremAnd project
\cite{PNTplus}, through the port in the \lean{zeta5} repository of C\'esar Del Solar \cite{zeta5}; the ported files
keep their Apache License~2.0. All error terms are explicit, except those coming from the prime number theorem and
from the compactness arguments in Section~\ref{sec:nonvan} (the decay constants of Lemma~\ref{lem:decay} and the
amplitude bounds $A_0,A_1$ of Lemma~\ref{lem:laplace}).

\subsection{How the certificates are checked}\label{sec:lean-cert}
The arithmetic certificate and the saddle certificate are decidable statements about rational data, proved by
\lean{decide +kernel}, i.e.\ by evaluation in the Lean kernel. The data (event lists, polyline vertices, rational
enclosures) were produced by scripts that are not trusted: a wrong value would make the build fail, but could not
make a false statement provable. The generators are in the repository (\lean{scripts\slash{}certA\slash{}},
\lean{scripts\slash{}saddle\slash{}gen\_saddle\_Z7.py}).

\emph{The arithmetic certificate} (Proposition~\ref{prop:certA}). The specification of $e^*$ and of the event list
(built with \lean{List.mergeSort}) is what the soundness proofs use, but the kernel cannot evaluate it efficiently.
The proof therefore passes through equivalent forms, each with a general proof of equivalence: a closed form of $e^*$
that uses only natural-number arithmetic per block (the kernel has fast \lean{Nat} primitives but none for
\lean{Int}); a decomposition of the event list into $133$ windows $[1+3j,4+3j)$, each sorted in the kernel and
compared with literal data, together with a check that their concatenation is strictly increasing; and a
decomposition of the rounded sum into $1{,}029$ chunks, each checked by a separate \lean{decide +kernel}. The resulting
value $U$ agrees exactly with an independent exact-rational Python implementation, and with the Python
recomputation \lean{scripts\slash{}certA\slash{}z7\_model.py} included in the repository, which mirrors the Lean
definitions.

\emph{The saddle certificate} (Theorem~\ref{thm:saddle}). This uses a verified interval arithmetic with rational
endpoints (rounded outwards to dyadic grids), complex boxes, verified enclosures of $\log$, $\arctan$, $\pi$ and
$\arg$, the contraction lemma for (S0), and segment, disc and ray checks as described in the proof of
Theorem~\ref{thm:saddle}. The phase multiplier ($2\pi$ for the $\beta$-analogue, $3\pi$ here) is a parameter of the
interval code.

\subsection{Cost}
With Mathlib cached, the arithmetic certificate dominates the build: its $1{,}029$ chunk lemmas live in $30$ files,
which take about $3.3$ CPU-hours in total and about $9.5$\,GB of memory each while the kernel evaluates them. The wall
time is therefore set by how many of these files fit in memory at once (\lean{LEAN\_NUM\_THREADS}): $2$ on a machine with
$32$\,GB of RAM, giving a build of a few hours, and $5$ with $64$\,GB. Everything else takes about $15$ minutes at $6$
threads. Running out of memory makes the build fail, but it is not a proof error, and rerunning the build continues
where it stopped.

\subsection{What has to be trusted}
Beyond the Lean kernel and the three axioms, the reader has to trust that the definition displayed above expresses
Theorem~\ref{thm:main}, and the Mathlib definitions it uses (\lean{riemannZeta}, real and complex numbers,
\lean{Complex.re}, \lean{Irrational}, \lean{Finset}). Nothing else, in particular no numerical data and no generating
script, needs to be trusted. Besides \lean{lake build}, the proof was replayed module by module with Lean's
\lean{leanchecker} ($134$ of the $141$ modules, including the main theorem and both certificates) and, in passes, with
the independent kernel \lean{nanoda} ($60{,}465$ of the $61{,}819$ theorems in the closure of the main theorem,
including the main theorem and every obligation of both certificates), with no error; these replays used the build of
an earlier commit that differs from the cited one only in comments and two lemma names (see \lean{TRUST.md} in the
repository). Apart from the check by the Lean kernel, all verification of the formalisation so far
was carried out by AI systems; no human expert has yet reviewed the Lean code.

% ===========================================================================
\section{Remarks}\label{sec:remarks}
% ===========================================================================

\subsection{The margin}
The proof needs $Q=\Re\Omega+\Acert<0$, and $Q\approx-0.198$. The two terms are of size about $10$, so the margin is
about $2\%$. Both numbers are certified: $\Acert$ by exact rational arithmetic and $\Re\Omega$ by interval arithmetic.
There is no third term: the $2$-adic and $3$-adic costs that twisted constructions pay are absent.

\subsection{How the parameters were found}\label{sec:search}
The parameters were found by a numerical search [numerical, not part of the proof]. A fast model of $Q$ (the
arithmetic integral on a coarse grid, and the saddle point or the real-axis exponent) was optimised over the block
profile, with offsets on the grid $\frac1{80}\ZZ\cap[0,\frac38]$, and over $r\in\frac1{24}\ZZ$ (later
$\frac1{48}\ZZ$), and the best candidates were then certified. In every run the optimum was close to a uniform
spread: one block at each of $S$ consecutive grid points $k/80$, with $r$ between about $0.8$ and $1$. Compared with
a single block at $\delta=0$, the spread profile is worth about $6$ per unit of $n$ at $S=27$, and the refinement
Corollary~\ref{cor:A}(b$'$) together with the exact counts of Section~\ref{sec:A} about $1$--$2$ per unit of $n$
[numerical].

For $S=17$ (which would also give $\zeta(7),\dots,\zeta(21)$) the best profile found has certified
$\Re\Omega+\Acert\approx+0.553$, so it does not close; with $S=18$ blocks the uniform profile with $r=11/12$ gives
Z7.

\subsection{\texorpdfstring{Why not $\zeta(19)$}{Why not zeta(19)}}\label{sec:N19}
To reach $\zeta(7),\dots,\zeta(19)$ this class needs $S\le16$. With uniform profiles the model gives $Q\approx+1.14$
for $S=16$ ($r=5/6$) and $Q\approx+1.82$ for $S=15$, including the saddle gain; a finer offset grid $\frac1{160}\ZZ$
at $S=17$ improved the model by less than $0.1$ [numerical]. So $N=19$ appears out of reach for this class
[heuristic]; the window $\zeta(7),\zeta(9),\zeta(11),\zeta(13)$ is about $4.5$ per unit of $n$ away.

One idea seemed to close the gap and fails. Replace the outer part of each zero run by zeros of lower multiplicity
$m_2$ (a \emph{tapered} run): multiplicity $5$ on $[1,r_1n]$ and $m_2$ on $(r_1n,r_2n]$, mirrored. Since only the
integers up to the crossing point $nx_c$ need quintuple zeros, the model suggested $Q\approx-0.8$ at $S=16$ when the
sum $\sum_tR_n''''(t)$ over the taper region was assumed to telescope. For $m_2\ge2$ it does not: at the integers of
the taper region $R_n''''$ equals the real-axis envelope times terms of type $\sin^2(\pi t)=(1-\cos2\pi t)/2$, which
alias to frequency $0$ under integer sampling, so that sum has the size of the real-axis maximum, not of the saddle.
The extended certificate confirmed this: the required bound on the taper region fails by $6.32$ and $4.37$ per unit of
$n$ for the two best candidates. For $m_2=1$ the constraints fail as well ($Q\approx+0.95$ at best) [numerical].
Ideas not tried include other zero multiplicities on the whole run, blocks of multiplicity $2$ at small $\delta$, and
genuinely different families.

\subsection{\texorpdfstring{Comparison with Lai--Zhou}{Comparison with Lai-Zhou}}
No $N$ makes ``one of $\zeta(7),\dots,\zeta(N)$'' imply ``two of $\zeta(5),\dots,\zeta(35)$'': the irrationality of
$\zeta(7)$ alone would satisfy the former. Conversely, Lai and Zhou's theorem gives only $N=35$ for windows starting at
$\zeta(7)$. If the irrationality of $\zeta(5)$ is accepted, Theorem~\ref{thm:main} is the stronger of the two
statements about $\zeta(7),\zeta(9),\dots$. A ``two of'' statement for windows starting at $\zeta(7)$ would need an
additional elimination, as in \cite{LaiZhou2022,RivoalZudilin2020}, which this family does not do. For the
construction of Lai and Zhou with simple zeros and twisted sums, the same search gave only $Q\approx-0.68$ at their
own window ($S=33$) and $Q\approx+0.37$ at $S=31$ in our model [numerical], so the derivative family is the better one
here.

\subsection{Effectivity}
The theorem is proved for $n\ge n_0$ with an unspecified $n_0$, because the prime number theorem is used without an
error term and some constants come from compactness. This does not affect the irrationality statement. The explicit
error terms are large: for instance the small-prime term $182\sqrt\cN\log\cN$ of Proposition~\ref{prop:decomp} falls
below $0.198\,n$ only for $n$ of order $2\cdot10^9$ [numerical estimate]. An effective version would need an explicit
form of the prime number theorem and of the compactness constants.

\subsection{What remains outside Lean}
Nothing that the proof of Theorem~\ref{thm:main} uses. The Lean statement is exactly the one displayed in
Section~\ref{sec:lean}, about Mathlib's \lean{riemannZeta}, and it uses no hypothesis on $\zeta(3)$ or $\zeta(5)$.
Outside Lean are only: the numerical search that found the parameters, the generators of the certificate data (both
untrusted and not needed for correctness), the exact checks at small $n$ quoted in Section~\ref{sec:general}, and the
certificates of the instances Z1--Z6, which are not used for Theorem~\ref{thm:main}.

\subsection{Review}
Before the formalisation, a prose version of the proof of Theorem~\ref{thm:general} (with the kernel split of
Remark~\ref{rem:split}) was reviewed adversarially by four independent AI reviewers (algebra, arithmetic, analysis, certificates and logic), who tested its claims with their
own code. They found no fatal or major mathematical issue; one presentation error (the values of $Q$ had been rounded
to nearest instead of up) and about twenty minor points were corrected. Among their checks: Lemma~A and Lemma~B in
$186$ small instances at every prime in range; the integrality of $p^{E_p}(\rho_n,\tau_i)$ at every prime for $16$
instances with $n\le32$; an independent evaluation of $e^*$ that reproduces the certified body sums; the identity
behind Remark~\ref{rem:split} at Z7 for $n=240$ and $480$, to $10^{-16}$; and a re-run of both certificates for
Z1, Z6 and Z7 [numerical]. The formalisation then found no gap; the points where the formal proof is organised
differently are those of Section~\ref{sec:lean-diff}.

\section*{Acknowledgements and disclosure}
The construction, the proofs, the parameter search, the certificate programs, the Lean formalisation and the text of
this paper were produced largely by AI models (Claude, Anthropic), working under the direction of the author. Apart
from the check by the Lean kernel, all verification so far, including the reviews mentioned above, was carried out by
AI systems. No human expert has yet reviewed the mathematics or the Lean code; we would be grateful for corrections.
The prime number theorem in the formalisation is ported from the PrimeNumberTheoremAnd project of A.~Kontorovich,
T.~Tao and collaborators, through the \lean{zeta5} repository of C\'esar Del Solar.

\bibliographystyle{amsplain}
\bibliography{references}

\end{document}
